% Options for packages loaded elsewhere
\PassOptionsToPackage{unicode}{hyperref}
\PassOptionsToPackage{hyphens}{url}
\PassOptionsToPackage{dvipsnames,svgnames,x11names}{xcolor}
%
\documentclass[
  11pt,
]{article}
\usepackage{amsmath,amssymb}
\usepackage{iftex}
\ifPDFTeX
  \usepackage[T1]{fontenc}
  \usepackage[utf8]{inputenc}
  \usepackage{textcomp} % provide euro and other symbols
\else % if luatex or xetex
  \usepackage{unicode-math} % this also loads fontspec
  \defaultfontfeatures{Scale=MatchLowercase}
  \defaultfontfeatures[\rmfamily]{Ligatures=TeX,Scale=1}
\fi
\usepackage{lmodern}
\ifPDFTeX\else
  % xetex/luatex font selection
\fi
% Use upquote if available, for straight quotes in verbatim environments
\IfFileExists{upquote.sty}{\usepackage{upquote}}{}
\IfFileExists{microtype.sty}{% use microtype if available
  \usepackage[]{microtype}
  \UseMicrotypeSet[protrusion]{basicmath} % disable protrusion for tt fonts
}{}
\makeatletter
\@ifundefined{KOMAClassName}{% if non-KOMA class
  \IfFileExists{parskip.sty}{%
    \usepackage{parskip}
  }{% else
    \setlength{\parindent}{0pt}
    \setlength{\parskip}{6pt plus 2pt minus 1pt}}
}{% if KOMA class
  \KOMAoptions{parskip=half}}
\makeatother
\usepackage{xcolor}
\usepackage[margin=2.2cm]{geometry}
\setlength{\emergencystretch}{3em} % prevent overfull lines
\providecommand{\tightlist}{%
  \setlength{\itemsep}{0pt}\setlength{\parskip}{0pt}}
\setcounter{secnumdepth}{5}
\usepackage{lmodern}
\usepackage[T1]{fontenc}
\usepackage{amsmath,amssymb}
\usepackage{newunicodechar}
\newunicodechar{}{\ensuremath{\bar\nu}}
\newunicodechar{}{\ensuremath{\hat\nu}}
\newunicodechar{}{\ensuremath{\tilde\sigma}}
\newunicodechar{}{\ensuremath{\bar v}}
\newunicodechar{}{\ensuremath{\hat\sigma}}
\newunicodechar{}{\ensuremath{\bar X}}
\newunicodechar{}{\ensuremath{\hat X}}
\newunicodechar{}{\ensuremath{\tilde W}}
\newunicodechar{}{\ensuremath{\dot\alpha}}
\newunicodechar{}{\ensuremath{\dot\beta}}
\newunicodechar{}{\ensuremath{\hat f}}
\newunicodechar{}{\ensuremath{\bar t}}
\newunicodechar{}{\ensuremath{\hat V}}
\newunicodechar{Ĉ}{\ensuremath{\hat C}}
\newunicodechar{ĉ}{\ensuremath{\hat c}}
\newunicodechar{ġ}{\ensuremath{\dot g}}
\newunicodechar{ᵀ}{\ensuremath{{}^{\mathsf{T}}}}
\newunicodechar{ℝ}{\ensuremath{\mathbb{R}}}
\newunicodechar{−}{\ensuremath{-}}
\newunicodechar{²}{\ensuremath{{}^{2}}}
\newunicodechar{³}{\ensuremath{{}^{3}}}
\newunicodechar{¹}{\ensuremath{{}^{1}}}
\newunicodechar{⁴}{\ensuremath{{}^{4}}}
\newunicodechar{⁸}{\ensuremath{{}^{8}}}
\newunicodechar{⁺}{\ensuremath{{}^{+}}}
\newunicodechar{⁻}{\ensuremath{{}^{-}}}
\newunicodechar{₂}{\ensuremath{{}_{2}}}
\newunicodechar{½}{\ensuremath{\tfrac12}}
\newunicodechar{¼}{\ensuremath{\tfrac14}}
\newunicodechar{·}{\ensuremath{\cdot}}
\newunicodechar{×}{\ensuremath{\times}}
\newunicodechar{÷}{\ensuremath{\div}}
\newunicodechar{±}{\ensuremath{\pm}}
\newunicodechar{σ}{\ensuremath{\sigma}}
\newunicodechar{φ}{\ensuremath{\varphi}}
\newunicodechar{ν}{\ensuremath{\nu}}
\newunicodechar{τ}{\ensuremath{\tau}}
\newunicodechar{π}{\ensuremath{\pi}}
\newunicodechar{ρ}{\ensuremath{\rho}}
\newunicodechar{λ}{\ensuremath{\lambda}}
\newunicodechar{μ}{\ensuremath{\mu}}
\newunicodechar{κ}{\ensuremath{\kappa}}
\newunicodechar{ω}{\ensuremath{\omega}}
\newunicodechar{β}{\ensuremath{\beta}}
\newunicodechar{α}{\ensuremath{\alpha}}
\newunicodechar{γ}{\ensuremath{\gamma}}
\newunicodechar{θ}{\ensuremath{\theta}}
\newunicodechar{η}{\ensuremath{\eta}}
\newunicodechar{χ}{\ensuremath{\chi}}
\newunicodechar{ψ}{\ensuremath{\psi}}
\newunicodechar{ξ}{\ensuremath{\xi}}
\newunicodechar{δ}{\ensuremath{\delta}}
\newunicodechar{ε}{\ensuremath{\varepsilon}}
\newunicodechar{Σ}{\ensuremath{\Sigma}}
\newunicodechar{Δ}{\ensuremath{\Delta}}
\newunicodechar{Λ}{\ensuremath{\Lambda}}
\newunicodechar{Π}{\ensuremath{\Pi}}
\newunicodechar{Ξ}{\ensuremath{\Xi}}
\newunicodechar{Γ}{\ensuremath{\Gamma}}
\newunicodechar{Φ}{\ensuremath{\Phi}}
\newunicodechar{Θ}{\ensuremath{\Theta}}
\newunicodechar{Ψ}{\ensuremath{\Psi}}
\newunicodechar{∫}{\ensuremath{\int}}
\newunicodechar{∂}{\ensuremath{\partial}}
\newunicodechar{√}{\ensuremath{\surd}}
\newunicodechar{∞}{\ensuremath{\infty}}
\newunicodechar{≈}{\ensuremath{\approx}}
\newunicodechar{≥}{\ensuremath{\ge}}
\newunicodechar{≤}{\ensuremath{\le}}
\newunicodechar{≠}{\ensuremath{\ne}}
\newunicodechar{≡}{\ensuremath{\equiv}}
\newunicodechar{→}{\ensuremath{\to}}
\newunicodechar{←}{\ensuremath{\leftarrow}}
\newunicodechar{⇒}{\ensuremath{\Rightarrow}}
\newunicodechar{⇔}{\ensuremath{\Leftrightarrow}}
\newunicodechar{↦}{\ensuremath{\mapsto}}
\newunicodechar{∈}{\ensuremath{\in}}
\newunicodechar{∝}{\ensuremath{\propto}}
\newunicodechar{∏}{\ensuremath{\prod}}
\newunicodechar{⊥}{\ensuremath{\perp}}
\newunicodechar{‖}{\ensuremath{\|}}
\newunicodechar{⟨}{\ensuremath{\langle}}
\newunicodechar{⟩}{\ensuremath{\rangle}}
\newunicodechar{∩}{\ensuremath{\cap}}
\newunicodechar{⊆}{\ensuremath{\subseteq}}
\newunicodechar{…}{\ensuremath{\ldots}}
\newunicodechar{̄}{}
\newunicodechar{̂}{}
\newunicodechar{̃}{}
\newunicodechar{̇}{}
\ifLuaTeX
  \usepackage{selnolig}  % disable illegal ligatures
\fi
\IfFileExists{bookmark.sty}{\usepackage{bookmark}}{\usepackage{hyperref}}
\IfFileExists{xurl.sty}{\usepackage{xurl}}{} % add URL line breaks if available
\urlstyle{same}
\hypersetup{
  pdftitle={Theory Question Bank},
  pdfauthor={Computational Finance (WI4151), TU Delft --- Prof.~Fang Fang},
  colorlinks=true,
  linkcolor={blue},
  filecolor={Maroon},
  citecolor={Blue},
  urlcolor={Blue},
  pdfcreator={LaTeX via pandoc}}

\title{Theory Question Bank}
\author{Computational Finance (WI4151), TU Delft --- Prof.~Fang Fang}
\date{}

\begin{document}
\maketitle

{
\hypersetup{linkcolor=blue}
\setcounter{tocdepth}{3}
\tableofcontents
}
Filled in per cluster as I study. Format: question -\textgreater{} my
short answer -\textgreater{} (gap?).

\hypertarget{cluster-1-foundations-lattice-pde}{%
\section{Cluster 1 --- Foundations / lattice /
PDE}\label{cluster-1-foundations-lattice-pde}}

\hypertarget{lecture-01}{%
\subsection{Lecture 01}\label{lecture-01}}

\begin{enumerate}
\def\labelenumi{\arabic{enumi}.}
\tightlist
\item
  Why does a course built around the COS method open with equity
  options, given that equities are a tiny fraction of total notional
  outstanding? \textbf{A.} Because the European equity call/put is the
  canonical optionality problem, and every later extension --- Heston,
  jumps, barriers, American/Bermudan --- is built on exactly the same
  scaffolding (risk-neutral discounted expectation of a payoff).
  Notional is dominated by rates, but the \emph{pricing theory} of
  optionality is cleanest and most transferable starting from a single
  stock under GBM, which is also where COS is first derived.
\item
  Derive the continuously-compounded bank-account value M(t) from the
  differential argument dM = rM dt, and state the discount factor for
  one euro at time T. \textbf{A.} Interest credited is proportional to
  balance, rate and time: dM = rM dt, so dM/M = r dt; integrating from 0
  to t gives ln M(t) − ln M(0) = rt, hence M(t) = M(0)e\^{}\{rt\}. With
  M(0)=1 the numéraire is e\^{}\{rt\}, and one euro at T is worth D(t,T)
  = e\^{}\{-r(T-t)\} at time t.
\item
  State covered interest-rate parity (1+i\_s) = (F\_t/S\_t)(1+i\_c) and
  prove it by a static no-arbitrage replication. What is the ``FX
  basis'' and why is it needed in practice? \textbf{A.} One domestic
  unit, two riskless routes over the same horizon: (A) invest
  domestically → 1+i\_c; (B) convert at spot S\_t, invest abroad at
  i\_s, lock reconversion at forward F\_t → (F\_t/S\_t)(1+i\_s).
  No-arbitrage forces equality, giving (1+i\_s)=(F\_t/S\_t)(1+i\_c) ---
  i.e.~the forward is pinned by spot and the two funding rates,
  model-free. Empirically CIP fails slightly; the FX basis is the spread
  added to one leg to restore equality, treated as calibrated market
  data fed into OTC FX pricing.
\item
  Write the terminal payoffs of a European call and put. Where are they
  non-smooth, and why does that kink cause numerical trouble for
  Fourier/COS methods? \textbf{A.} V\_C = max(S\_T−K,0), V\_P =
  max(K−S\_T,0): continuous, piecewise-linear, with a kink at S=K.
  Expanding the \emph{payoff} itself in a Fourier/cosine series rings
  (Gibbs) at the kink, so the cosine coefficients decay only
  algebraically; COS sidesteps this by expanding the smooth
  \emph{density} and integrating the payoff against cosines analytically
  (the χ\_k, ψ\_k closed forms), so the kink never enters a series.
\item
  In the one-step binomial hedge of a short call, derive Delta =
  20/(S\_up - S\_dn) and then the premium V\_\{c,0\}. Compute both for
  \{120,80\} and \{120,75\}. Why is case B more expensive? \textbf{A.}
  Hedged portfolio Π = V\_\{c,0\} − ΔS\_0; set Π\_up = Π\_dn so ΔS\_up −
  20 = ΔS\_dn (with K=100, up-payoff 20), giving Δ = 20/(S\_up−S\_dn).
  Imposing Π=0 gives V\_\{c,0\} = Δ(S\_0−S\_dn). Case A \{120,80\}: Δ=½,
  V=½·20=10. Case B \{120,75\}: Δ=4/9, V=4/9·25=11.11. Case B is dearer
  because the wider outcome spread (more volatility) is worth more to a
  convex, floored payoff.
\item
  Show that the one-step replication price equals E\^{}Q{[}payoff{]}
  under the risk-neutral measure, and that the real-world probabilities
  p=0.7, p=0.4 never enter. Define q explicitly. \textbf{A.} Define q =
  (S\_0−S\_dn)/(S\_up−S\_dn) so that
  E\^{}Q{[}S\_T{]}=qS\_up+(1−q)S\_dn=S\_0 (martingale under Q, r=0).
  Then V\_\{c,0\}=Δ(S\_0−S\_dn)=q(S\_up−K)+(1−q)·0 =
  E\^{}Q{[}(S\_T−K)\^{}+{]}. The price is built purely from the
  no-arbitrage hedge; the physical p (0.7, 0.4) never appears --- only q
  does.
\item
  At the fair price, does a perfectly hedged writer make money? Explain
  Pi\_up = Pi\_dn = 0 and where real-world profit actually comes from.
  \textbf{A.} No.~At the replication price the perfectly hedged terminal
  portfolio is zero in \emph{both} states (Π\_up=Π\_dn=0): no risk, no
  profit --- the fair price is exactly the cost of replication. Profit
  only arises from selling above replication value (the surplus invested
  at r), and in reality from bid--ask spread, fees,
  discrete-rehedging/transaction-cost edge, and bearing vol/jump risk
  outside the model.
\item
  Derive put-call parity from two portfolios held to maturity. Which
  assumptions about the dynamics of S did you use? \textbf{A.} Π\_A =
  call + K cash, Π\_B = put + share. At T both equal max(S\_T,K) (using
  max(S\_T−K,0)+K and max(K−S\_T,0)+S\_T). Equal payoffs in every state
  ⇒ equal value today: C + Ke\^{}\{-r(T-t)\} = P + S\_t. No dynamics
  assumed --- only no-arbitrage and a constant rate --- so parity is
  model-free.
\item
  Prove the lower bound C(t) \textgreater= max(S(t) - K
  e\^{}\{-r(T-t)\}, 0) and use it to argue an American call on a
  non-dividend stock is never exercised early. Why does the argument
  fail for puts? \textbf{A.} The portfolio (long call + bond worth
  Ke\^{}\{-r(T-t)\}) pays max(S\_T,K)≥S\_T at T, so it costs ≥ S(t)
  today: C(t) ≥ S(t)−Ke\^{}\{-r(T-t)\}, and ≥0, hence the bound. For
  r\textgreater0 this exceeds intrinsic S(t)−K, so the live call beats
  immediate exercise and American = European. For the put, parity gives
  P\^{}eur = Ke\textsuperscript{\{-r(T-t)\}−S+C}eur and
  Ke\^{}\{-r(T-t)\}\textless K, so the live put can be \emph{below}
  intrinsic K−S --- banking K early to earn interest can win, so early
  exercise is genuinely optimal.
\item
  Apply Ito's lemma to g = ln S under GBM and obtain d(ln S) = (mu -
  sigma\^{}2/2) dt + sigma dW. Where does the -1/2 sigma\^{}2 come from,
  and what goes wrong if you forget it? \textbf{A.} With g=ln S,
  g\_S=1/S, g\_SS=−1/S²: d ln S = (μS·1/S + ½σ²S²·(−1/S²))dt +
  σS·(1/S)dW = (μ−½σ²)dt + σdW. The −½σ² is the Itô correction, born of
  (dW)²=dt in the second-order term. Forgetting it biases the
  simulated/expected stock high: it is exactly the convexity adjustment
  that keeps E{[}S(t){]}=S\_0e\^{}\{μt\}.
\item
  From the log-price ABM, write the closed-form S(t) and verify
  E{[}S(t){]} = S\_0 e\^{}\{mu t\} using the Gaussian MGF. State the
  distribution of ln(S(T)/S\_0) under Q. \textbf{A.} Integrating d ln S
  gives S(t)=S\_0 exp{[}(μ−½σ²)t + σW(t){]}. Using
  E{[}e\^{}\{σW(t)\}{]}=e\^{}\{σ²t/2\}: E{[}S(t){]}=S\_0
  e\textsuperscript{\{(μ−σ²/2)t\}·e}\{σ²t/2\}=S\_0 e\^{}\{μt\}. Under Q
  (μ→r), ln(S(T)/S\_0) \textasciitilde{} N((r−½σ²)T, σ²T) --- lognormal
  stock.
\item
  State the four defining properties of a Wiener process. What are
  E{[}dW{]}, E{[}(dW)\^{}2{]}, and the quadratic variation {[}W{]}\_t,
  and why does that force Ito calculus instead of ordinary calculus?
  \textbf{A.} (i) W(0)=0; (ii) independent, stationary increments; (iii)
  W(t)\textasciitilde N(0,t); (iv) a.s. continuous paths. So
  E{[}dW{]}=0, E{[}(dW)²{]}=dt, and {[}W{]}\_t=t (finite quadratic
  variation, unbounded total variation). Because (dW)²=dt contributes at
  first order in dt, second-order Taylor terms survive --- ordinary
  calculus fails and Itô's correction ½σ²g\_xx dt is needed.
\item
  Define the Markov property via P(S\_\{t+1\} \textless= s \textbar{}
  S\_t) = P(S\_\{t+1\} \textless= s \textbar{} S\_t, S\_\{t-1\},
  \ldots). How does it relate to the weak-form EMH, and why does i.i.d.
  returns imply weak-form EMH but not conversely? \textbf{A.} Markov:
  the one-step-ahead conditional law given the present equals that given
  the whole past --- history beyond S\_t carries no extra information,
  which is precisely weak-form efficiency (prices reflect past prices).
  i.i.d. returns make the \emph{entire} history useless, which certainly
  gives Markov/weak-form; but weak-form only requires the one-step
  conditional law to be unchanged by deeper history, a strictly weaker
  condition than full independence and identical distribution --- so the
  converse fails.
\item
  Contrast ABM, GBM and OU: which can go negative, which is
  multiplicative, which mean-reverts, and what single property do all
  three share? Which one is the template for the Heston variance
  process? \textbf{A.} ABM (dX=μdt+σdW): additive noise, Gaussian, can
  go negative. GBM (dX=μXdt+σXdW): multiplicative, stays positive,
  lognormal. OU (dX=κ(θ−X)dt+σdW): mean-reverts to θ at speed κ,
  Gaussian, can go negative. Shared property: all are Markov diffusions
  driven by Brownian motion. OU is the template for the CIR/Heston
  variance dv=κ(θ−v)dt+γ√v dW --- same mean-reverting drift but a √v
  diffusion that keeps v≥0.
\item
  What stylised fact does the SPX QQ-plot reveal (sample quantiles
  {[}-3,3{]} vs theoretical {[}-2.3,2.3{]}), and which later models in
  the course are introduced to fix it? \textbf{A.} Empirical return
  quantiles (≈±3) are wider than the normal's (≈±2.3): returns have
  \emph{fatter tails} than Gaussian (large moves occur more often). This
  is the central failure of GBM and motivates Heston stochastic
  volatility (L7) and jump-diffusion models (Merton/Kou, L8), which put
  weight back in the tails and reproduce the implied-vol smile.
\end{enumerate}

\hypertarget{lecture-02}{%
\subsection{Lecture 02}\label{lecture-02}}

\begin{enumerate}
\def\labelenumi{\arabic{enumi}.}
\tightlist
\item
  Derive Ito's lemma heuristically: why does the second-order term in
  g(t,X) survive while the others vanish? \textbf{A.} Taylor-expand dg =
  g\_t dt + g\_x dX + ½g\_xx(dX)² + \ldots{} With dX=μdt+σdW,
  (dX)²=σ²(dW)²+2μσ dt dW+μ²(dt)². Under the table (dW)²=dt, dW dt=0,
  (dt)²=0, only σ²(dW)²=σ²dt survives at order dt. So
  dY=(g\_t+μg\_x+½σ²g\_xx)dt+σg\_x dW: the second-order term survives
  precisely because Brownian motion has finite quadratic variation
  ((dW)²=dt), while genuinely higher-order terms vanish.
\item
  Apply Ito to g = ln S under dS = mu S dt + sigma S dW. Where does the
  -sigma\^{}2/2 come from, and what is its sign relation to Jensen?
  \textbf{A.} g\_S=1/S, g\_SS=−1/S²: d ln S = (μ−½σ²)dt+σdW. The −½σ² is
  the Itô term from ½g\_SS(dS)² = −½σ²dt. It is Jensen's inequality made
  precise: ln is concave, so E{[}ln S\_t{]} grows at the lower rate
  μ−½σ² while E{[}S\_t{]} grows at μ; the gap is exactly ½σ².
\item
  Walk through the delta-hedging derivation of the Black-Scholes PDE. At
  which exact step is Delta = V\_S forced, and why? \textbf{A.} Form
  Π=V−ΔS, self-financing so dΠ=dV−ΔdS. Itô gives
  dV=(V\_t+μSV\_S+½σ²S²V\_SS)dt+σSV\_S dW. Then
  dΠ=(\ldots)dt+(σSV\_S−ΔσS)dW. The portfolio is locally riskless iff
  the dW term vanishes, which forces Δ=V\_S. Imposing dΠ=rΠdt then
  yields V\_t+rSV\_S+½σ²S²V\_SS−rV=0.
\item
  In that derivation, mu appears in dV but not in the final PDE. Show
  precisely the cancellation that removes it. \textbf{A.} In dΠ the
  drift is V\_t+μSV\_S+½σ²S²V\_SS−Δ·μS. Once Δ=V\_S is substituted, the
  term μSV\_S (from dV) and −Δ·μS=−μSV\_S (from −ΔdS) cancel exactly,
  removing μ entirely. Economically: hedging eliminates exposure to the
  real-world drift, so only σ and r remain.
\item
  State the BS PDE, its order/type, and its terminal condition. Why is
  it solved backward in time? \textbf{A.} V\_t+rSV\_S+½σ²S²V\_SS−rV=0:
  second-order, linear, parabolic, with terminal data V(S,T)=payoff
  g(S). It is solved backward because the only data we know is at
  maturity (the payoff); discounted expectation propagates value from T
  back to t, equivalently the substitution τ=T−t turns it into a forward
  initial-value problem.
\item
  Define a risk-neutral measure Q. Starting from dS = mu' S dt + sigma S
  dW\^{}Q, prove that the martingale property of e\^{}\{-rt\}S\_t forces
  mu' = r. \textbf{A.} Q is a measure equivalent to P under which the
  discounted asset e\^{}\{-rt\}S\_t is a martingale. By the product rule
  d(e\^{}\{-rt\}S\_t) = (μ'−r)e\^{}\{-rt\}S\_t dt + σe\^{}\{-rt\}S\_t
  dW\^{}Q. The stochastic integral is a martingale; the drift
  dt-integral is a martingale only if it vanishes, i.e.~μ'−r=0. Hence
  μ'=r and dS=rSdt+σS dW\^{}Q.
\item
  Under the change from P to Q, the drift goes mu -\textgreater{} r but
  sigma is unchanged. Why must sigma be preserved (think equivalence /
  quadratic variation)? \textbf{A.} Girsanov shifts only the drift; the
  diffusion coefficient is invariant. Equivalence of P and Q means they
  agree on which events have probability zero, and σ is fixed by the
  path's quadratic variation {[}S{]}\_t=∫σ²S²ds --- an almost-sure
  pathwise property. Changing σ would change the set of possible paths'
  quadratic variation, breaking equivalence; so only the drift can move.
\item
  State Feynman-Kac and sketch its proof. What is the role of the
  discounted process Y\_s and where is the PDE used? \textbf{A.} If V
  solves V\_t+bV\_x+½σ²V\_xx−rV=0 with V(x,T)=g, and dX=b dt+σ dW with
  X\_t=x, then V(x,t)=E{[}e\^{}\{−∫\_t\^{}T r
  du\}g(X\_T)\textbar X\_t=x{]}. Proof: set Y\_s=e\^{}\{−∫\_t\^{}s r
  du\}V(X\_s,s); Itô gives dY\_s =
  (discount)·{[}V\_t+bV\_x+½σ²V\_xx−rV{]}ds + (\ldots)dW. The PDE makes
  the ds-drift vanish, so Y\_s is a martingale; taking E of Y\_T−Y\_t
  with Y\_t=V(x,t) gives the formula.
\item
  Explain how Feynman-Kac shows the PDE approach and the martingale
  approach give the same price. State the equivalence ``drift = 0
  \textless=\textgreater{} discounted price is a Q-martingale''.
  \textbf{A.} The discounted option price e\^{}\{-rt\}V is a
  Q-martingale iff its Itô drift e\^{}\{-rt\}(V\_t+rSV\_S+½σ²S²V\_SS−rV)
  is zero --- which is exactly the BS PDE. So ``PDE holds'' ⇔
  ``discounted price is a Q-martingale'', and applying the martingale
  property between t and T gives
  V=e\textsuperscript{\{-r(T-t)\}E}Q{[}g(S\_T){]}. PDE route and
  expectation route are two faces of the same identity.
\item
  Derive C\_0 = e\^{}\{-rT\} E\^{}Q{[}(S\_T - K)\^{}+{]} down to S\_0
  N(d1) - K e\^{}\{-rT\} N(d2). What measure-change / square-completion
  produces the shift from d2 to d1? \textbf{A.} Write
  C\_0=e\^{}\{-rT\}∫\_\{ln(K/S\_0)\}\^{}∞(S\_0 e\^{}z−K)f\_Z(z)dz with
  Z\textasciitilde N((r−½σ²)T,σ²T). The K-term is
  Ke\textsuperscript{\{-rT\}·Q(S\_T\textgreater K)=Ke}\{-rT\}N(d\_2).
  For the S\_0 e\^{}z-term, completing the square in z−(z−m)²/2v² gives
  S\_0 e\^{}z f\_Z(z)=S\_0 e\^{}\{rT\}·(N(m+v²,v²) density), i.e.~a
  Girsanov shift of the mean up by v²=σ²T (change to the stock
  numéraire). That shift turns d\_2 into d\_1=d\_2+σ√T, giving
  C\_0=S\_0N(d\_1)−Ke\^{}\{-rT\}N(d\_2).
\item
  Interpret N(d2) and N(d1) probabilistically. Which one equals Q(S\_T
  \textgreater{} K)? \textbf{A.} N(d\_2)=Q(S\_T\textgreater K), the
  risk-neutral probability of finishing in the money. N(d\_1) is the
  same exercise probability but under the \emph{stock-numéraire} measure
  Q\^{}S; equivalently S\_0N(d\_1) is the present value (in numéraire
  terms) of receiving the stock conditional on exercise. So
  Q(S\_T\textgreater K)=N(d\_2).
\item
  Derive put-call parity from a static replication argument (not from
  the BS formula). Why is parity model-free? \textbf{A.} Long one call,
  short one put (same K,T) pays
  (S\_T−K)\textsuperscript{+−(K−S\_T)}+=S\_T−K in every state. A payoff
  S\_T−K at T is replicated today by one share and short K bonds, value
  S\_0−Ke\^{}\{-rT\}. No-arbitrage ⇒ C\_0−P\_0=S\_0−Ke\^{}\{-rT\}. It is
  model-free because it uses only the algebraic payoff identity, a
  static portfolio, and a constant rate --- no assumption on the
  dynamics of S.
\item
  Greeks: define Delta, Gamma, Theta. Show the BS PDE as an algebraic
  identity linking them, and give the closed-form call Delta and Gamma.
  \textbf{A.} Δ=∂V/∂S, Γ=∂²V/∂S², Θ=∂V/∂t. The BS PDE is literally
  Θ+rSΔ+½σ²S²Γ=rV --- an algebraic relation tying the Greeks together
  (they are not independent). For the call Δ=N(d\_1)∈(0,1) and
  Γ=N'(d\_1)/(S\_0σ√T)\textgreater0.
\item
  List the five parameters the BS call price depends on and explain why
  mu (the real-world drift) is not among them. \textbf{A.} The call
  depends on S, K, T (or τ), r, σ. The physical drift μ is absent
  because delta-hedging removes exposure to it (the μSV\_S terms
  cancel), equivalently because under Q the drift is replaced by r. Risk
  preferences are hedged away; only the cost of carry r and the
  volatility σ matter.
\end{enumerate}

\hypertarget{lecture-03}{%
\subsection{Lecture 03}\label{lecture-03}}

\begin{enumerate}
\def\labelenumi{\arabic{enumi}.}
\tightlist
\item
  In the implied-volatility problem, which four inputs of the BS price
  are known and which one is the unknown? Why is recovering σ a
  root-find rather than a closed-form inversion? \textbf{A.} Known: τ,
  K, r, S\_t (all observable). Unknown: σ. Recovering σ from C\^{}mkt is
  a root-find of F(σ)=C\textsuperscript{BS(σ)−C}mkt=0 because σ sits
  inside the nonlinear N(d\_1), N(d\_2); there is no algebraic inversion
  of the normal CDF, so we iterate.
\item
  Prove vega = ∂C/∂σ = S0√T φ(d1) \textgreater{} 0 for σ\textgreater0.
  Why does strict monotonicity of price in σ guarantee a unique implied
  vol? \textbf{A.} Differentiating C=S\_tN(d\_1)−Ke\^{}\{-rτ\}N(d\_2)
  and using the identity S\_tN'(d\_1)=Ke\^{}\{-rτ\}N'(d\_2) (from
  N'(d\_1)/N'(d\_2)=e\^{}\{−ln(S/K)−rτ\} since d\_2=d\_1−σ√τ), the
  ∂d\_1/∂σ terms cancel and Vega = Ke\^{}\{-rτ\}N'(d\_2)√τ = S\_t√τ
  N'(d\_1) \textgreater{} 0. Strict positivity ⇒ C(σ) is strictly
  increasing, hence a continuous bijection onto its range, so at most
  one σ solves F(σ)=0.
\item
  State the admissible no-arbitrage price interval for a European call.
  What happens to the implied vol as the market price approaches the
  upper/lower bound? \textbf{A.} As σ→0⁺, C→max(S\_t−Ke\^{}\{-rτ\},0)
  (intrinsic); as σ→∞, C→S\_t. So C is a bijection from {[}0,∞) onto
  {[}max(S\_t−Ke\^{}\{-rτ\},0), S\_t). A quote at the lower bound forces
  σ\^{}impl→0; a quote approaching the upper bound S\_t forces
  σ\^{}impl→∞. A quote outside the band has no real implied vol
  (static-arbitrage violation / stale quote).
\item
  Bisection vs Newton-Raphson for IV inversion: state the convergence
  order of each and Newton's failure mode. Why is dividing by vega
  dangerous deep OTM / short maturity? \textbf{A.} Bisection: linear,
  rate ½ (one bit/step), globally fail-safe given a sign-changing
  bracket. Newton σ\_\{n+1\}=σ\_n−F(σ\_n)/Vega(σ\_n): quadratic but only
  local --- a poor start can overshoot, diverge, or step to negative σ.
  Vega=S\_t√τN'(d\_1) is a Gaussian bump that →0 deep ITM/OTM and at
  short τ; dividing by a near-zero vega makes the Newton step explode,
  so the inversion is ill-conditioned in the wings.
\item
  Globalised Newton: why seed at the price-curve inflection point ?
  What property there makes Newton converge monotonically (no
  overshoot)? \textbf{A.} ∂²C/∂σ² ∝ d\_1d\_2, which vanishes at
  =√(2\textbar ln(S\_t/K)+rτ\textbar/τ): C is convex for σ\textless,
  concave for σ\textgreater, so  is the inflection point and where
  Vega=C' is \emph{maximal}. Starting there, the contraction ratio
  1−F'(ξ)/F'()∈(0,1) because F'() is the maximum slope, so the first
  step cannot overshoot the root; the iterates then stay on the monotone
  (concave) side and converge monotonically and quadratically.
\item
  Put implied vol: why must a put and a call at the same (K,T) have the
  same implied vol? Derive it from put-call parity. \textbf{A.} Parity
  C−P=S\_t−Ke\^{}\{-rτ\} is σ-independent, so
  P(σ)=C(σ)−(S\_t−Ke\^{}\{-rτ\}). Hence ∂P/∂σ=∂C/∂σ=Vega\textgreater0:
  the put has the same vega and strict monotonicity. Since the market
  call and put at (K,T) satisfy parity, the σ that reprices the call
  also reprices the put --- one implied vol per (K,T).
\item
  Describe the implied-vol smile/skew in the strike direction and the
  term structure in maturity. What market fact does each encode?
  \textbf{A.} Strike direction: a \emph{smile} (convex U, high vol in
  both tails) encodes excess kurtosis / heavy tails vs lognormal; a
  \emph{skew/smirk} (decreasing in K, high vol for low-strike puts)
  encodes negative skewness / crash risk, the typical equity-index
  shape. Maturity direction (term structure): flat under constant σ;
  under time-dependent σ(t), σ\^{}impl(T)²=(1/T)∫\_0\^{}T σ(t)²dt (RMS
  average), so a rising instantaneous vol gives an upward-sloping term
  structure. Generic surface: pronounced short-maturity smile flattening
  to a milder long-maturity skew.
\item
  Why does constant-vol Black-Scholes fail to fit the whole surface?
  State precisely what a single σ cannot reproduce. \textbf{A.} If BS
  held, σ\^{}impl would be one constant for all (K,T) --- a flat plane.
  The market surface is curved and tilted, so a single σ cannot
  reproduce simultaneously the strike-direction smile/skew (heavy,
  asymmetric tails) and the maturity term structure (vol clustering /
  time-varying variance). Constant volatility and Gaussian returns are
  the broken assumptions.
\item
  Distinguish implied vs local vs stochastic volatility: which is a
  quoting convention, which is a deterministic σ(S,t), which adds its
  own SDE? How does this motivate Heston (Lecture 7)? \textbf{A.}
  Implied vol is a model-inversion \emph{output} --- a re-quoting of
  price in BS units, one number per quote, saying nothing about true
  dynamics. Local vol σ(S,t) is a deterministic \emph{input} surface
  (Dupire) calibrated to fit the whole surface today. Stochastic vol
  σ(W\^{}v,t) is a \emph{random} input with its own driving SDE (and
  correlation ρ with spot). Implied describes the data, local fits it
  deterministically, stochastic explains it --- Heston (random
  mean-reverting variance, L7) generates smile/skew endogenously and is
  the COS workhorse.
\end{enumerate}

\hypertarget{lecture-05}{%
\subsection{Lecture 05}\label{lecture-05}}

\begin{enumerate}
\def\labelenumi{\arabic{enumi}.}
\tightlist
\item
  Derive the CRR up/down/probability parameters by matching the mean and
  variance of the tree log-return to Black-Scholes. Show every step.
  \textbf{A.} One step: ln(S/S\_0)=ln d + ln(u/d)·R,
  R\textasciitilde Bernoulli(p). Mean = p ln u+(1−p)ln d, Var =
  p(1−p)(ln(u/d))². Match to BS log-return over δt: mean=(r−½σ²)δt,
  var=σ²δt. So p ln u+(1−p)ln d=(r−½σ²)δt and ln(u/d)=σ√δt/√(p(1−p)).
  Closing with p=½ gives ln u−ln d=2σ√δt and ln u+ln d=2(r−½σ²)δt, hence
  u=exp(σ√δt+(r−½σ²)δt), d=exp(−σ√δt+(r−½σ²)δt).
\item
  Why does the Itô-corrected drift r-½σ² (not r) appear inside u and d?
  What exactly is being matched -- the price or the log-price?
  \textbf{A.} The matching is on the \emph{log}-price increment, whose
  risk-neutral mean is (r−½σ²)δt by Itô (d ln S=(r−½σ²)dt+σdW). So the
  drift slot in u,d absorbs the Itô-corrected drift r−½σ², not r;
  matching Var(ln S)=σ²δt reproduces the diffusion coefficient.
\item
  The moment equations give two equations in three unknowns (p,u,d).
  Which free choice does the deck make, and how does that differ from
  the classical d=1/u CRR tree? Do both match the first two moments?
  \textbf{A.} Two equations, three unknowns ⇒ one free choice. The deck
  picks p=½ (symmetric), giving the u,d above. Classical CRR instead
  imposes d=1/u and solves for p, giving u=e\^{}\{σ√δt\},
  d=e\^{}\{−σ√δt\}, p=(e\^{}\{rδt\}−d)/(u−d). Both match the first two
  moments to O(δt); they differ only in the higher-order/free-parameter
  choice.
\item
  State and justify the convergence rate of the binomial method. Why is
  the convergence non-monotone (sign-oscillating in M)? What is the
  mechanism? \textbf{A.} First order, \textbar error\textbar≤K/M, but
  non-monotone (sign oscillates in M). Mechanism: the discrete terminal
  nodes S\_n\^{}M straddle the strike K differently as M changes; the
  position of K relative to the lattice wobbles irregularly, so the
  payoff-region sampling error oscillates. This is why a raw tree cannot
  be naively Richardson-extrapolated.
\item
  Write the FTCS stencil for u\_t=u\_xx, the explicit update, and the
  matrix form U\^{}\{i+1\}=F U\textsuperscript{i+p}i. Why do boundary
  terms appear only in the first and last entries of p\^{}i? \textbf{A.}
  Forward time, central space:
  (U\_j\textsuperscript{\{i+1\}−U\_j}i)/k=(U\_\{j+1\}\textsuperscript{i−2U\_j}i+U\_\{j−1\}\^{}i)/h²,
  so
  U\_j\textsuperscript{\{i+1\}=νU\_\{j+1\}}i+(1−2ν)U\_j\textsuperscript{i+νU\_\{j−1\}}i
  with ν=k/h². In matrix form U\textsuperscript{\{i+1\}=FU}i+p\^{}i,
  F=tridiag(ν,1−2ν,ν), p\^{}i=(νa, 0,\ldots,0, νb)\^{}T. Boundary terms
  enter only the first and last entries because only the stencils at j=1
  and j=N\_x−1 reach into the known boundary nodes U\_0, U\_\{N\_x\}.
\item
  Carry out the von Neumann analysis for FTCS. Derive ξ(β)=1-4ν
  sin²(βh/2) and show the stability condition reduces to ν=k/h²≤½. Which
  mode is worst? \textbf{A.} Insert U\_j\textsuperscript{i=ξ}i
  e\^{}\{iβjh\}: ξ=1+ν(e\textsuperscript{\{iβh\}−2+e}\{−iβh\})=1−4ν
  sin²(βh/2), using e\textsuperscript{\{iθ\}−2+e}\{−iθ\}=−4sin²(θ/2).
  ξ≤1 automatically; the binding constraint ξ≥−1 gives 4ν sin²(βh/2)≤2.
  The worst mode is the highest frequency βh=π (sin²=1), giving ν≤½,
  i.e.~k≤h²/2.
\item
  Do the same for BTCS: derive ξ=1/(1+4ν sin²(βh/2)) and explain why
  BTCS is unconditionally stable. \textbf{A.} Backward time evaluates
  the space operator at the new level: ξ−1=−4νξ sin²(βh/2), so ξ=1/(1+4ν
  sin²(βh/2)) ∈ (0,1{]} for every ν\textgreater0. Since
  \textbar ξ\textbar≤1 for all modes and all k, there is no step
  restriction --- BTCS is unconditionally stable. (CN, ξ=(1−2ν
  sin²)/(1+2ν sin²), is likewise unconditionally stable.)
\item
  What is the CFL-type cost of FTCS's explicitness in terms of how k
  must scale with h? Why is this expensive when refining the grid?
  \textbf{A.} Stability forces k≤h²/2, so the time step scales
  quadratically with the space step: halving h forces a fourfold cut in
  k (and so \textasciitilde4× the time steps). Refining space therefore
  makes explicit FTCS very expensive --- the cost of explicitness.
  Implicit BTCS/CN avoid this with one tridiagonal solve per step.
\item
  Give the local truncation errors of FTCS, BTCS and Crank-Nicolson. Why
  does CN gain an order in time while costing the same per step as BTCS?
  \textbf{A.} FTCS: O(k)+O(h²); BTCS: O(k)+O(h²) (time term opposite
  sign); CN: O(k²)+O(h²). CN is centred at the half time level i+½,
  which by symmetry kills the leading O(k) term, leaving O(k²). It is
  still one tridiagonal solve per step like BTCS, so it gains an order
  in time at the same per-step cost --- hence preferred (caveat: it can
  oscillate on non-smooth payoffs).
\item
  State the Lax equivalence theorem and explain its role: consistency +
  stability ⇒ convergence. Why isn't consistency alone enough?
  \textbf{A.} For a consistent linear finite-difference scheme of a
  well-posed linear IVP, the scheme converges iff it is stable.
  Consistency (truncation error →0) only says the scheme approximates
  the PDE locally; without stability, errors (especially high-frequency
  modes, \textbar ξ\textbar\textgreater1) are amplified each step and
  blow up --- so the local approximation does not accumulate to a
  convergent global solution. Stability bounds the error growth, and the
  two together give convergence.
\item
  Transform the reverse-time BS PDE to a constant-coefficient equation
  via X=ln S, then remove the reaction term via V=e\^{}\{-rτ\}W. Show
  the coefficient algebra (S²∂\_SS=∂\_XX-∂\_X). \textbf{A.} X=ln S gives
  ∂\_S=(1/S)∂\_X and ∂\_SS=(1/S²)(∂\_XX−∂\_X), so S²∂\_SS=∂\_XX−∂\_X and
  S∂\_S=∂\_X. The reverse-time PDE V\_τ−½σ²S²V\_SS−rSV\_S+rV=0 becomes
  V\_τ−½σ²V\_XX−(r−½σ²)V\_X+rV=0 (constant coefficients). Then
  V=e\^{}\{-rτ\}W: V\_τ=e\^{}\{-rτ\}(W\_τ−rW), and the rV and −rW terms
  cancel, leaving W\_τ−½σ²W\_XX−(r−½σ²)W\_X=0, a pure
  advection--diffusion equation.
\item
  Show that applying BTCS to the log-transformed advection-diffusion
  equation with the coupling h²=σ²k makes the centre node drop out,
  yielding
  V\textsuperscript{i\_j=e}\{-rk\}(p*V\^{}\{i+1\}\_\{j+1\}+(1-p*)V\textsuperscript{\{i+1\}\emph{\{j-1\}).
  Derive p\emph{. \textbf{A.} BTCS on W\_τ=½σ²W\_XX+(r−½σ²)W\_X produces
  the centre coefficient
  (1/k)W\_j\textsuperscript{\{i+1\}−(σ²/h²)W\_j}\{i+1\}. With the
  coupling h²=σ²k, σ²/h²=1/k, so the two centre terms cancel and only
  the j±1 neighbours survive:
  W\_j\textsuperscript{i=p\emph{W\_\{j+1\}\^{}\{i+1\}+(1−p})W\_\{j−1\}}\{i+1\}.
  Collecting the up-coefficient (diffusion 1/2k plus convection
  (r−½σ²)/2h) and using h=σ√k gives p}=½(1+(r−½σ²)√k/σ). Restoring
  V=e\^{}\{-rτ\}W multiplies one τ-step by e\^{}\{-rk\}:
  V\_j\textsuperscript{i=e}\{-rk\}(p\emph{V\_\{j+1\}\^{}\{i+1\}+(1−p})V}\{j−1\}}\{i+1\}).
\item
  In what precise sense is the binomial method ``an explicit FD scheme
  customised to a single time-zero value''? What does a general FD
  scheme give that the tree throws away? \textbf{A.} The recursion above
  is exactly the risk-neutral binomial backward step (up-prob p\emph{,
  log-spacing ±σ√k), i.e.~an explicit FD scheme on the transformed heat
  equation, ``customised'' via h²=σ²k so the centre node drops out and
  only the two branches remain. The tree returns only V\_0\^{}0 (one
  spot). A general FD scheme returns the value at }every* grid node
  \{jh,ik\}, hence prices for all initial spots at once --- that
  whole-grid information is what the tree discards.
\item
  On the raw S-grid the FTCS diffusion weight is ½kσ²j². What does this
  do to stability near S=L, and how does working in x=ln S fix it?
  Connect this to domain-truncation pitfalls (cf.~the COS {[}a,b{]}
  range at small c\_2). \textbf{A.} The weight ½kσ²j² grows like j², so
  the effective stability condition σ²j²k\textless\textasciitilde1 is
  hardest at the largest j (S near the truncation L): stability degrades
  quadratically near the upper boundary. Working in x=ln S makes the PDE
  constant-coefficient, so the stencil weight is uniform and this
  blow-up disappears. Both are the same lesson as the COS {[}a,b{]}
  trap: a method on a finite window is only as good as the window --- FD
  truncates the state space, COS truncates the Fourier-cosine range, and
  at small c\_2=σ²T that COS window collapses just as the FD domain must
  be chosen to hold all the mass.
\item
  The deck's Slide 57 prints a garbled p* formula. Why is it wrong, and
  what is the correct dimensionless form p\emph{=½(1+(r-σ²/2)√k/σ)?
  Check the k→0 limit. \textbf{A.} The slide prints ½(1+√k/σ·r−σ²),
  which is dimensionally inconsistent (the −σ² is not scaled and a
  factor is missing). The correct algebra gives p}=½(1+(r−½σ²)√k/σ),
  which is dimensionless (the √k/σ carries units that cancel). As k→0
  the correction →0 and p*→½, consistent with the symmetric p=½ choice
  of the moment-matched tree. Treat the slide as a transcription error.
\end{enumerate}

\hypertarget{cluster-2-monte-carlo}{%
\section{Cluster 2 --- Monte Carlo}\label{cluster-2-monte-carlo}}

\hypertarget{lecture-04}{%
\subsection{Lecture 04}\label{lecture-04}}

\begin{enumerate}
\def\labelenumi{\arabic{enumi}.}
\tightlist
\item
  Construct the MC integral estimator and prove it is unbiased. Derive
  Var(\_N)=σ²/N from independence and state the standard error.
  \textbf{A.} Write I=∫h g dx=E\_g{[}h(X){]}, estimate by
  \_N=(1/N)Σh(x\_i), x\_i\textasciitilde g i.i.d. Unbiased:
  E{[}\_N{]}=(1/N)ΣE{[}h(X){]}=I. Variance:
  Var(\_N)=(1/N²)ΣVar(X\_i)=(1/N²)·Nσ²=σ²/N (cross-covariances vanish
  by independence). Standard error se=σ/√N=O(N\^{}\{-1/2\}).
\item
  Why is the MC error O(N\^{}\{-1/2\}) independent of dimension, while
  tensor-product quadrature is O(N\^{}\{-k/d\})? When does MC win?
  \textbf{A.} The CLT rate se=σ/√N depends only on the number of samples
  N and the variance σ², not on the dimension d --- sampling a
  high-dimensional density is no harder per draw. Tensor-product
  quadrature of order k uses N=m\^{}d points (m per axis) for
  O(m\textsuperscript{\{−k\})=O(N}\{−k/d\}), which degrades as d grows
  (curse of dimensionality). MC wins in high dimension: baskets, many
  time steps, path-dependent payoffs.
\item
  Write a CLT-based confidence interval for the MC estimate. Why use the
  sample variance s\_N² with (N-1), and what does it cost? \textbf{A.}
  By the CLT (\_N−I)/(σ/√N)→N(0,1), so a (1−α) CI is \emph{N ±
  z}\{1−α/2\}·σ/√N (z=1.96 for 95\%). σ is unknown, so use
  s\_N²=(1/(N−1))Σ(X\_i−\_N)²; the N−1 (Bessel) makes it an
  \emph{unbiased} variance estimate, at the cost of one degree of
  freedom (negligible for large N). The reported half-width z·s\_N/√N is
  exactly what variance reduction tries to shrink.
\item
  State the Koksma-Hlawka inequality (star discrepancy + Hardy-Krause
  variation). Why can QMC achieve near-O(N\^{}\{-1\}) = (log N)\^{}d/N?
  \textbf{A.} \textbar(1/N)Σf(x\_i)−∫f\textbar{} ≤ V(f)·D\_N\^{}\emph{,
  where V(f) is the Hardy--Krause variation (integrand roughness,
  independent of points) and D\_N\^{}} is the star discrepancy
  (worst-case box gap, independent of f). A low-discrepancy sequence
  (Sobol, Halton) has D\_N\^{}*=O((log N)\^{}d/N), so for
  bounded-variation f the deterministic QMC error is O((log N)\^{}d/N)
  --- essentially O(N\^{}\{−1\}), a full power faster than plain MC's
  O(N\^{}\{−1/2\}).
\item
  What is ``effective dimension'' and why can plain QMC degrade in high
  nominal dimension? What does randomised QMC recover? \textbf{A.} The
  (log N)\^{}d factor is enormous for large d, so the QMC advantage only
  shows at impractical N unless the \emph{effective} dimension (how much
  each coordinate actually matters) is low --- Brownian-bridge / PCA
  path constructions lower it and make QMC work for option pricing.
  Plain QMC is deterministic, giving no error estimate; randomised QMC
  (e.g.~random shifts of a lattice rule) restores an unbiased estimator
  and lets you estimate the error from independent shifts while keeping
  most of the discrepancy gain.
\item
  Write the Euler-Maruyama scheme for dX=a dt+b dW. What iterated Itô
  integral does it drop? \textbf{A.}
  \emph{\{t\_k\}=}\{t\_\{k−1\}\}+a(t\_\{k−1\},)h+b(t\_\{k−1\},)√h
  Z\_k, Z\_k\textasciitilde N(0,1), freezing coefficients at the left
  endpoint (Itô convention). It drops the iterated Itô integral ∫∫dW
  dW=½{[}(ΔW)²−h{]}, i.e.~the variation of b \emph{within} the step ---
  the leading correction that Milstein restores.
\item
  Derive the Milstein correction from the Itô-Taylor expansion: show
  ∫∫dW dW = ½{[}(ΔW)²-h{]} and the resulting ½bb'{[}(ΔW)²-h{]} term.
  \textbf{A.} Expand σ(X\_u) by Itô around X\_s: dσ(X\_t)=\ldots dt+σ'σ
  dW, so σ(X\_u)=σ(x)+σ'(x)σ(x)∫dW+O(dt). Substituting into ∫σ dW gives
  the extra term σ'σ·∫\_s\^{}\{t\_k\}∫\_s\^{}u dW\_r dW\_u. The inner
  integral is ∫(W\_u−W\_s)dW\_u=½{[}(ΔW)²−h{]}. So the update gains
  ½σσ'{[}(ΔW)²−h{]}=½bb'(Z²−1)h. It vanishes for additive noise (b
  constant), where Euler already has strong order 1.
\item
  Define strong vs weak convergence order. Give the orders for Euler
  (½,1) and Milstein (1,1). Why does Milstein buy pathwise but not
  distributional accuracy? \textbf{A.} Strong order β:
  E\textbar\_\{Kh\}−X\_T\textbar≤ch\^{}β (pathwise, same Brownian
  path). Weak order β:
  \textbar E{[}f(\_\{Kh\}){]}−E{[}f(X\_T){]}\textbar≤ch\^{}β for smooth
  f (laws only). Euler: strong ½, weak 1; Milstein: strong 1, weak 1.
  Milstein improves the \emph{strong} (pathwise) rate but has the same
  weak rate as Euler --- so for a European price E{[}f(X\_T){]} it buys
  nothing distributionally; it matters for path-dependent payoffs and
  for coupling coarse/fine levels (multilevel MC) where pathwise
  accuracy is used.
\item
  Antithetic variates: derive the variance of the antithetic estimator
  and state the monotonicity/comonotonicity condition for it to help.
  \textbf{A.} Pair each draw U with its reflection 1−U:
  Î=(1/M)Σ½(f(U\_i)+f(1−U\_i)), unbiased. Per pair,
  Var(½(f(U)+f(1−U)))=½Var(f(U))+½Cov(f(U),f(1−U)). Comparing with the
  plain estimator on the same 2M evaluations, antithetic wins iff
  Cov(f(U),f(1−U))\textless0, which is guaranteed when f is monotone (by
  the FKG/correlation lemma: f increasing ⇒ −f(1−·) increasing ⇒
  Cov(f(U),f(1−U))≤0). It can backfire for strongly non-monotone payoffs
  (e.g.~a straddle).
\item
  Control variates: derive the optimal coefficient
  θ\emph{=Cov(X,Y)/Var(Y) and show the variance is reduced by (1-ρ²).
  What is a good control for an Asian option? \textbf{A.} With
  Z\_θ=X+θ(E{[}Y{]}−Y), E{[}Z\_θ{]}=E{[}X{]} for all θ (unbiased),
  Var(Z\_θ)=Var(X)−2θCov(X,Y)+θ²Var(Y). Minimising the quadratic gives
  θ}=Cov(X,Y)/Var(Y), and substituting,
  Var(Z\_\{θ\emph{\})=(1−ρ²\_\{XY\})Var(X). The reduction is dramatic
  when \textbar ρ\textbar{} is near 1. For an arithmetic-average Asian
  (no closed form), use the }geometric*-average Asian (closed form for
  equally spaced fixings) as control --- the two payoffs are nearly
  perfectly correlated path by path (Kemna--Vorst).
\item
  Importance sampling: write the likelihood-ratio estimator and the
  Gaussian drift-shift. What is the infinite-variance pitfall, and how
  does Girsanov connect? \textbf{A.} Sample from a proposal q and
  reweight: E\_p{[}f(X){]}=E\_q{[}f(X)p(X)/q(X){]}, estimator
  (1/N)Σf(X\_i)p(X\_i)/q(X\_i), X\_i\textasciitilde q, with p/q the
  likelihood ratio (dp/dq). Gaussian shift X=Z+θ:
  V=E{[}f(Z+θ)exp(−θᵀZ−½‖θ‖²){]}, moving mass into the payoff region
  (e.g.~so S\_T\textgreater K is typical for a deep-OTM call). Pitfall:
  a proposal with too-light tails makes p/q blow up and the estimator
  has \emph{infinite variance} --- worse than plain MC; IS is the only
  one of the four that can backfire catastrophically. In continuous time
  the reweighting is exactly Girsanov: shifting the drift of W by Y\_s
  corresponds to the density process exp(∫Y dW−½∫Y²ds).
\item
  Why is MC the benchmark that COS and FFT are compared against, and on
  what axes (rate, dimension, smoothness) does COS dominate for European
  options? \textbf{A.} MC is unbiased by construction and
  dimension-robust, so it is the brute-force reference: a fast method is
  ``right'' if it agrees with the MC price to within the CLT standard
  error. For a 1-D European under a model with known smooth
  characteristic function, COS converges \emph{exponentially} in the
  number of cosine terms, vs MC's O(N\^{}\{−1/2\}) and FFT's algebraic
  quadrature rate --- reaching machine precision with a few hundred
  terms where MC needs \textasciitilde10\^{}\{12\} paths. MC still wins
  on high-dimensional baskets and path-dependent payoffs with no
  transform structure.
\end{enumerate}

\hypertarget{cluster-3-fourier-cos}{%
\section{Cluster 3 --- Fourier / COS}\label{cluster-3-fourier-cos}}

\begin{quote}
COS-specific drilling lives in \texttt{cos-deep-dive.md}. This section
holds the surrounding Fourier-family theory from Lecture 6 that Fang can
ask ``around'' COS.
\end{quote}

\hypertarget{lecture-6-the-fourier-family-around-cos}{%
\subsection{Lecture 6 --- the Fourier family (around
COS)}\label{lecture-6-the-fourier-family-around-cos}}

\begin{enumerate}
\def\labelenumi{\arabic{enumi}.}
\tightlist
\item
  State the continuous Fourier transform and its inverse. Which
  scaling/sign convention did the slides use, and why does the
  convention not matter for the method? \textbf{A.} One convention:
  (ω)=∫f(x)e\^{}\{i2πωx\}dx, f(x)=(1/2π)∫(ω)e\^{}\{−i2πωx\}dω. The
  slides place the 2π and sign so that the characteristic function φ\_X
  is (up to that 2π-scaling) the Fourier transform of the density. The
  convention does not matter for COS because only internal consistency
  is needed --- COS reads cosine coefficients off φ at the grid
  frequencies k π/(b−a), and any consistent choice gives the same A\_k.
\item
  Define the characteristic function φ\_X(t). In what precise sense is
  it the ``Fourier dual'' of the density f\_X? Write the inversion
  formula f\_X(x) = (1/2π)∫ e\^{}\{-itx\} φ\_X(t) dt. \textbf{A.}
  φ\_X(t)=E{[}e\^{}\{itX\}{]}=∫e\^{}\{itx\}f\_X(x)dx --- this \emph{is}
  the Fourier transform of the density (up to the 2π convention), so
  f\_X and φ\_X are a Fourier pair. The inversion is
  f\_X(x)=(1/2π)∫e\^{}\{−itx\}φ\_X(t)dt. This is exactly the integral
  COS evaluates cheaply, without quadrature.
\item
  Why are characteristic-function methods attractive at all? Name four
  models whose ch.f. is known in closed form but whose density is not.
  \textbf{A.} Because for a huge class of models the density f is
  intractable but the characteristic function φ is known in closed (or
  semi-closed) form, and a method needing only φ inherits that whole
  library for free. Examples: Heston, Variance-Gamma, Merton/Kou
  jump-diffusions (and the affine jump-diffusions generally), CGMY/NIG
  Lévy models.
\item
  Write the DFT and inverse DFT. What does the FFT change about the
  cost, from what to what, and by what factorisation idea? Does COS use
  the FFT? \textbf{A.} DFT X\_k=Σ\_n x\_n e\^{}\{i2πkn/N\}, inverse
  x\_n=(1/N)Σ\_k X\_k e\^{}\{−i2πkn/N\}. Naively O(N²); the FFT
  factorises the DFT matrix into O(log N) sparse factors (radix-2 split
  into even/odd sub-DFTs of length N/2), giving O(N log N). European COS
  does \emph{not} use the FFT --- its cost is the bare O(N) evaluation
  of one cosine sum (the FFT returns only for the
  discrete-barrier/Bermudan recursions).
\item
  Derive the half-range cosine series from the full Fourier series via
  the even extension. Why does the even extension kill the sine terms?
  \textbf{A.} A function on {[}0,L{]} has a full Fourier series with sin
  and cos terms. Extending it \emph{evenly} to {[}−L,L{]} makes it an
  even function; since sin is odd, all sine coefficients b\_n=(2/L)∫f
  sin vanish, leaving the half-range cosine series f(x)=a\_0/2+Σ a\_n
  cos(nπx/L), a\_n=(2/L)∫\_0\^{}L f cos(nπx/L)dx. COS uses this on
  {[}a,b{]} (shift x↦x−a, period 2(b−a)); the cosine basis gives real
  coefficients read directly off φ.
\item
  Carr--Madan: why is the undamped call C\_T(k) not square-integrable in
  log-strike k? What is the role of the damping factor e\^{}\{αk\}, and
  what is the cost of having a free parameter α? \textbf{A.} As k=ln
  K→−∞, C\_T(k)→S\_0 (a nonzero constant), so C\_T is not in L¹/L² and
  cannot be Fourier-transformed directly. Damping
  c\_T(k)=e\^{}\{αk\}C\_T(k) with α\textgreater0 makes it integrable so
  its Fourier transform exists. The cost is the free parameter α: it has
  no canonical value, its choice affects accuracy/stability, and it must
  be tuned. COS needs no such parameter (the range comes from
  cumulants).
\item
  Carr--Madan: ψ\_T(v) = e\^{}\{-rT\} φ\_X(v-(α+1)i) /
  (α²+α-v²+i(2α+1)v). Why is it a selling point that ψ\_T is expressible
  through the ch.f.? How is the price recovered, and by what numerical
  tool? \textbf{A.} Because the transform of the damped call ψ\_T is
  written purely through φ\_X (evaluated at the shifted complex argument
  v−(α+1)i), the whole library of models with known φ is immediately
  available --- no density needed. The price is recovered by the inverse
  transform C\_T(k)=(e\textsuperscript{\{−αk\}/2π)∫e}\{−ivk\}ψ\_T(v)dv,
  evaluated numerically by the FFT (a trapezoidal-type quadrature of an
  oscillatory integrand). The two weaknesses vs COS: it still does a
  numerical \emph{integration}, and it needs α.
\item
  Contrast COS and Carr--Madan on: core operation, free parameters,
  convergence rate for smooth f, and number of terms to reach machine
  precision. \textbf{A.} COS: one O(N) cosine sum, no integration;
  \emph{no} free parameter (range from cumulants); exponential
  convergence for smooth f; dozens--hundreds of terms to machine
  precision; analytic Greeks from the same sum. Carr--Madan/FFT: an FFT
  quadrature O(N log N); free damping α to tune; algebraic
  (quadrature-limited) convergence; thousands of nodes for far worse
  accuracy. Both need only φ.
\item
  Give the BS log-price characteristic function and state the
  risk-neutral drift μ = r - ½σ². Where does the martingale condition
  fix μ? \textbf{A.} φ(u)=exp(iuμT−½σ²u²T) with μ=r−½σ² for
  X\_T=ln(S\_T/S\_0). The drift is fixed by the martingale condition
  that e\^{}\{−rt\}S\_t be a Q-martingale: with S\_T=S\_0 exp(μT+σW\_T),
  E{[}e\^{}\{−rT\}S\_T{]}=S\_0 requires
  e\textsuperscript{\{μT\}E{[}e\^{}\{σW\_T\}{]}=e}\{rT\}, i.e.~μ+½σ²=r,
  so μ=r−½σ². (This is the same −½σ² Itô correction as in d ln S.)
\end{enumerate}

\hypertarget{lecture-10-barrier-options-theory-around-cos}{%
\subsection{Lecture 10 --- barrier options (theory around
COS)}\label{lecture-10-barrier-options-theory-around-cos}}

\begin{quote}
Deep COS-for-barrier drilling is in \texttt{cos-deep-dive.md} (sets
H--K). These are the surrounding theory items.
\end{quote}

\begin{enumerate}
\def\labelenumi{\arabic{enumi}.}
\setcounter{enumi}{9}
\tightlist
\item
  Name the 8 barrier types and explain why exactly 8. State in--out
  parity and its uses. \textbf{A.} up/down × in/out × call/put = 2×2×2 =
  8: up-and-out call, up-and-in call, down-and-out call, down-and-in
  call, and the four put analogues. In--out parity: knock-in + knock-out
  (same K,B,T) = vanilla, because exactly one of the two is alive at
  maturity --- C\_in(t)+C\_out(t)=e\^{}\{−r(T−t)\}E{[}(S\_T−K)\^{}+{]}.
  Uses: price the easier one (usually the knock-out) plus the vanilla
  and back out the other by subtraction; and a free correctness check on
  any numerical barrier price.
\item
  Write the barrier price as (i) a first-passage expectation and (ii) a
  localized Feynman--Kac PDE. How does the barrier enter the PDE?
  \textbf{A.} (i)
  V(x,t)=e\^{}\{−r(T−t)\}E\_\{x\}{[}(α(e\textsuperscript{\{X\_T\}−K))}+·1\_\{τ\_B\textgreater T\}{]}
  --- discounted vanilla payoff times a survival indicator. (ii) On
  {[}a,b{]}, V solves the localized BS PDE
  V\_t+(r−½σ²)V\_x+½σ²V\_xx−rV=0 with terminal data the vanilla payoff
  and \emph{absorbing (Dirichlet) boundary} V=0 at the knock-out
  level(s). The barrier enters not through the equation but through the
  homogeneous boundary condition: hitting the barrier sets value to
  zero.
\item
  State and prove the reflection principle result P(τ\_a ≤ T) = 2P(W\_T
  ≥ a). Hence write the density of the running maximum. \textbf{A.} By
  the strong Markov property at τ\_a and BM symmetry, a path that has
  touched a is equally likely to end above or below a:
  P(W\_T≥a\textbar τ\_a≤T)=½. So
  P(τ\_a≤T)=P(τ\_a≤T,W\_T≥a)+P(τ\_a≤T,W\_T\textless a)=2P(τ\_a≤T,W\_T≥a)=2P(W\_T≥a)
  (since reaching W\_T≥a forces τ\_a≤T). Hence the running max Ξ=max W
  has the law of \textbar W\_T\textbar, with density
  p\_Ξ(a)=√(2/πT)·e\^{}\{−a²/2T\}, a≥0.
\item
  Derive the joint density of (running max, W\_T). Why does the
  reflection at level 2a−b appear? \textbf{A.} Reflect the path after
  τ\_a across level a: for b≤a, P(Ξ≥a, W\_T\textless b)=P(Ξ≥a,
  W\_T\textgreater2a−b)=P(W\_T\textgreater2a−b) (since 2a−b≥a forces the
  max to exceed a). Differentiating the tail
  P(Ξ\textgreater a,W\_T≤b)=P(W\_T≥2a−b) twice gives
  p\_\{Ξ,W\_T\}(a,b)=√(2/π)·((2a−b)/T\textsuperscript{\{3/2\})·e}\{−(2a−b)²/2T\},
  a≥max(b,0). The 2a−b is the mirror image of the endpoint b reflected
  across the barrier a --- the reflected path's terminal point. (Drift
  is handled by a Girsanov reweighting.)
\item
  Method of images: state U(S,t) = S\^{}\{2α\}V(B²/S,t), α = ½(1−2r/σ²),
  and the down-and-out call it gives. \textbf{A.} One checks directly
  that if LV=0 then the image U(S,t)=S\^{}\{2α\}V(B²/S,t), α=½(1−2r/σ²),
  also solves LV=0 --- the BS analogue of reflecting a heat-equation
  solution across a wall. Subtracting it with the right weight cancels
  the value on S=B, giving the down-and-out call
  C\_DO(S\_t,K)=C\_van(S\_t,K)−(S\_t/B)\^{}\{2α\}C\_van(B²/S\_t,K): the
  vanilla minus the image evaluated at the mirror spot B²/S\_t with
  weight (S\_t/B)\^{}\{2α\}, which removes exactly the barrier-crossing
  paths.
\item
  Sine-series PDE route: what substitution turns the localized BS PDE
  into the heat equation, and why the sine (not cosine) basis?
  \textbf{A.} With τ=½σ²(T−t), κ=2r/σ², V=e\^{}\{αx+βτ\}U, α=−½(κ−1),
  β=−α²−κ, z=x−a, the localized BS PDE becomes the heat equation
  U\_τ=U\_zz on {[}0,b−a{]} with U(0,τ)=U(b−a,τ)=0. The homogeneous
  \emph{Dirichlet} (zero) boundary conditions select the \emph{sine}
  basis (sin(nπz/(b−a)) vanishes at both ends, cos does not); each mode
  propagates by e\^{}\{−(nπ/(b−a))²τ\}. This is the structural precursor
  of COS --- but the propagator is the explicit heat semigroup, whereas
  COS uses φ and so generalises to any model.
\item
  Why is a discretely-monitored barrier a backward recursion? Where does
  the Markov property enter, and why does the FFT reappear (vs European
  COS)? \textbf{A.} Survival is a product of indicators
  ∏1\_\{x\_m\textless h\}, and by the Markov property the joint density
  factorises into one-step transitions
  f(x\_M\textbar x\_\{M−1\})\ldots f(x\_1\textbar x\_0), so the
  M-dimensional integral \emph{nests} into a backward recursion: at each
  date c(x,t\_\{m−1\})=e\^{}\{−rΔt\}∫v(y,t\_m)f(y\textbar x)dy
  (European-COS form), then knock out by overwriting v=R above h, which
  splits the coefficient V\_k at h. The split feeds the previous V\_j
  through an analytic coupling M\_\{k,j\} that is Hankel-plus-Toeplitz,
  i.e.~a discrete convolution --- so each step is done by FFT in O(N log
  N), overall O((M−1)N log₂N). European COS needs no FFT (single O(N)
  sum); the barrier needs the whole coefficient vector propagated every
  step, which is the convolution the FFT accelerates.
\end{enumerate}

\hypertarget{cluster-4-models-heston-jumps}{%
\section{Cluster 4 --- Models (Heston /
jumps)}\label{cluster-4-models-heston-jumps}}

\hypertarget{lecture-07}{%
\subsection{Lecture 07}\label{lecture-07}}

\begin{enumerate}
\def\labelenumi{\arabic{enumi}.}
\tightlist
\item
  Write down the Heston SDE system under Q, naming every parameter, and
  state what dW\^{}S dW\^{}v = rho dt means physically. \textbf{A.}
  dS\_t=rS\_t dt+√ν\_t S\_t dW\^{}S\_t, dν\_t=κ(−ν\_t)dt+γ√ν\_t
  dW\^{}ν\_t, with dW\^{}S dW\^{}ν=ρ dt. Here ν\_t is the instantaneous
  \emph{variance} (vol is √ν\_t),  the long-run variance,
  κ\textgreater0 the mean-reversion speed, γ\textgreater0 the vol-of-vol
  (smile curvature), ρ the spot/variance correlation (skew slope).
  dW\^{}S dW\^{}ν=ρ dt means shocks to price and to variance are
  instantaneously correlated; ρ\textless0 (the leverage effect) means
  down-moves in S come with up-moves in variance, steepening the skew.
\item
  Derive the log-spot form d ln S\_t and show how the Cholesky factor
  decorrelates the two Brownian motions into independent W\^{}S, W\^{}v.
  \textbf{A.} Itô on ln S with (dS/S)²=ν dt: d ln S\_t=(r−½ν\_t)dt+√ν\_t
  dW\^{}S\_t. For Σ={[}{[}1,ρ{]},{[}ρ,1{]}{]}=LLᵀ,
  L={[}{[}1,0{]},{[}ρ,√(1−ρ²){]}{]}, write the correlated pair via
  independent (\_1,\_2): W\^{}S=\_1, W\^{}ν=ρ\_1+√(1−ρ²)\_2, which
  checks Cov(W\textsuperscript{S,W}ν)=ρt. Equivalently, driving the
  variance by its own clean W\^{}ν, the log-spot diffusion becomes
  √ν\_t(ρ dW\textsuperscript{ν+√(1−ρ²)dW}⊥) with W\textsuperscript{ν⊥W}⊥
  --- the form used for Euler and the affine derivation.
\item
  State the conditional law of nu\_T given nu\_t (noncentral
  chi-squared): give d, the noncentrality lambda, and the scaling c.
  \textbf{A.} ν\_T\textbar ν\_t \textasciitilde{} Y/(2c) with
  Y\textasciitilde χ*²\_d(λ) noncentral chi-squared,
  c=2κ/(γ²(1−e\^{}\{−κτ\})), τ=T−t, degrees of freedom d=4κ/γ²,
  noncentrality λ=2c e\^{}\{−κτ\}ν\_t. So the variance is exactly a
  scaled noncentral χ² --- this is what the Broadie--Kaya exact scheme
  samples.
\item
  State the Feller condition precisely. What exactly does it guarantee,
  and what happens to nu\_t when it is violated? \textbf{A.} Feller: 2κ
  ≥ γ². It guarantees the origin is inaccessible --- ν\_t stays strictly
  positive a.s. If violated, ν\_t can reach 0 (and reflects/leaves again
  for the pure CIR process), but it never goes negative. Practically the
  exact noncentral-χ² scheme is guaranteed well-defined (d\textgreater0)
  under Feller.
\item
  Derive the Feller condition from the noncentral chi-squared density
  blowing up at zero (use I\_nu(y) \textasciitilde{} (y/2)\^{}nu /
  Gamma(nu+1) as x-\textgreater0). \textbf{A.} The noncentral χ² density
  ∝ x\^{}\{d/4−1/2\}I\_\{d/2−1\}(√(λx)). As x→0,
  I\_\{d/2−1\}(√(λx))\textasciitilde(√(λx)/2)\^{}\{d/2−1\}/Γ(d/2), so
  the density behaves like
  x\textsuperscript{\{d/4−1/2\}·x}\{(d/2−1)/2\}=x\^{}\{d/2−1\}. For it
  not to blow up at the origin we need d/2−1≥0, i.e.~d=4κ/γ²≥2, which
  is 2κ≥γ² --- the Feller condition. A large γ relative to the
  restoring force κ makes excursions to zero likely.
\item
  Is the Heston model affine in (S, nu)? Show why not, then show how
  moving to ln S restores the affine condition (give K0, K1, H0, H1).
  \textbf{A.} Not affine in (S,ν): the diffusion covariance has top-left
  entry νS², not affine in the state. Switching to x=ln S, the drift
  becomes K0=(r, κ)ᵀ + K1·(x,ν)ᵀ with K1={[}{[}0,−½{]},{[}0,−κ{]}{]}
  (affine), and σσᵀ=ν·{[}{[}1,γρ{]},{[}γρ,γ²{]}{]}, i.e.~H0=0 and H1 the
  tensor whose ν-slice is {[}{[}1,γρ{]},{[}γρ,γ²{]}{]} and x-slice is 0
  --- affine in ν. So in log-coordinates Heston is affine and the
  Duffie--Pan--Singleton transform applies.
\item
  Set up the affine ansatz phi = exp(alpha + beta·x) and derive the
  Riccati ODE system. Where does each ODE come from when you substitute
  into the Feynman-Kac PDE? \textbf{A.} Feynman--Kac: φ solves
  ∂\emph{tφ+Lφ−Rφ=0, φ(T,x)=e\^{}\{u·x\}. Substituting φ=e\^{}\{α+β·x\}
  gives ∂}\{x\_i\}φ=β\_iφ, ∂\_\{x\_ix\_j\}φ=β\_iβ\_jφ, ∂\_tφ=(+·x)φ.
  Dividing by φ and using that μ, σσᵀ, R are affine in x, collect terms:
  the part \emph{linear} in x must vanish, giving =ρ\_1−K\_1ᵀβ−½βᵀH\_1β
  (the Riccati ODE); the part \emph{constant} in x gives
  =ρ\_0−K\_0·β−½βᵀH\_0β. Boundary conditions β(T,T)=u, α(T,T)=0.
  Affineness is what lets the x-dependence separate cleanly.
\item
  Solve the scalar CIR Riccati dot-beta = -kappa beta + (1/2) gamma\^{}2
  beta\^{}2 explicitly via g = 1/beta. Then integrate for alpha.
  \textbf{A.} Set g=1/β; then ġ=−/β²=κg−½γ², a linear ODE. With
  integrating factor e\^{}\{−κt\}: d(g
  e\textsuperscript{\{−κt\})=−½γ²e}\{−κt\}dt, so
  g(t)=e\textsuperscript{\{κt\}(g(0)−γ²(1−e}\{−κt\})/(2κ)), g(0)=1/u.
  Inverting, β(t)=u e\^{}\{−κt\}/(1−½uγ²ξ\_κ(t)) with
  ξ\_κ(t)=(1−e\^{}\{−κt\})/κ. Then =κβ integrates (since
  e\^{}\{−κs\}=ξ\_κ') to α(t)=−(2κ/γ²)ln(1−½uγ²ξ\_κ(t)). This is the
  same Riccati mechanism behind the full Heston cf.
\item
  Write the closed-form Heston cf of ln S\_T: give b, a, s and the
  formulas for beta(tau,u) and alpha(tau,u). \textbf{A.} With b=γρu−κ,
  a=u(1−u), s=√(b²+aγ²):
  β(τ,u)=a(1−e\textsuperscript{\{−sτ\})/(2s−(s+b)(1−e}\{−sτ\})),
  α(τ,u)=−rτ+ruτ−κ{[}(s+b)τ/γ²+(2/γ²)ln(1−((s+b)/2s)(1−e\^{}\{−sτ\})){]}.
  Then φ\_\{ln S\_T\}(ω;τ)=exp(α(τ,iω)+iω ln S\_t+β(τ,iω)ν\_t). This is
  the load-bearing object COS consumes.
\item
  Explain the ``little Heston trap'': what goes wrong at long
  maturities, what causes it (branch cut of complex log / sqrt), and
  which grouping fixes it? \textbf{A.} The cf involves a complex √ (s)
  and complex log, which are multi-valued. The original Heston (1993)
  grouping (with g=(b+s)/(b−s) and e\^{}\{+sτ\}) repeatedly crosses the
  principal branch cut of ln as τ grows, making the integrand
  discontinuous and giving wildly wrong long-maturity prices. The
  algebraically equivalent ``little trap'' grouping (Albrecher et
  al.~2007) --- written with e\^{}\{−sτ\} and
  1−((s+b)/2s)(1−e\^{}\{−sτ\}), the form above, taking the principal √
  with non-negative real part --- stays on the principal branch and is
  stable for all τ.
\item
  How does the Heston cf enter the COS pricing formula? At which
  frequencies is it evaluated, and why is Heston the ideal use case for
  COS? \textbf{A.} The Heston φ is the \emph{only} model-specific input
  to the COS formula v≈e\^{}\{−rτ\}Σ'
  Re\{φ(kπ/(b−a))e\^{}\{−ikπa/(b−a)\}\}V\_k, evaluated at the COS
  frequencies ω\_k=kπ/(b−a) (the V\_k are the model-free payoff
  coefficients). Heston is the ideal use case because its \emph{density}
  is unknown in closed form but its cf is known closed-form --- exactly
  the regime where COS (which never touches the density) dominates,
  reaching \textasciitilde10\^{}\{−10\} error with N≈160 terms.
\item
  Under Heston, how do you choose the COS truncation range {[}a,b{]}?
  Why is the short-maturity (small c2) case dangerous, and how does
  negative rho make it worse? \textbf{A.} Same cumulant rule
  {[}a,b{]}=c\_1±L√(c\_2+√c\_4), with c\_1,c\_2,c\_4 of ln S\_T obtained
  by differentiating log φ at ω=0. Danger at short τ: c\_2 (the
  variance) is small, so the box collapses toward c\_1; if the higher
  cumulants are under-weighted the range chops off mass and COS loses
  spectral accuracy --- adding terms N does not help (it is a range, not
  a series, error). Negative ρ produces a heavy left tail (skew) plus γ
  fattens the tails, so the true support is wider than a c\_2-only box;
  the fix is a larger L (10--12) and keeping the √c\_4 term, not the
  BS-style L=8--10 with c\_4 ignored. This is the Assignment-2 bug
  generalised.
\item
  Why can the Euler scheme produce negative variance? List truncation,
  reflection, and log-variance fixes and state the drawback of each.
  \textbf{A.} Euler updates ν\_\{i+1\}=ν\_i+κ(−ν\_i)Δ+γ√ν\_i √Δ Z with
  Z Gaussian and unbounded below, so
  P(ν\_\{i+1\}\textless0)\textgreater0 always; then √ν is undefined.
  Fixes: (i) truncation/floor ν←max(ν,ε) --- simple but introduces
  truncation bias; (ii) reflection ν←\textbar ν\textbar{} --- useful
  when Feller fails but the reflected path sits systematically above the
  floored one near zero; (iii) simulate ln ν by Euler and exponentiate
  --- does \emph{not} work, ln ν can blow up so ν=e\^{}\{lnν\} explodes.
  All are biased, worst when Feller is badly violated.
\item
  State the Broadie-Kaya exact simulation algorithm (4 steps). Which
  step is hard and how is it solved (cf of integrated variance
  -\textgreater{} COS inversion -\textgreater{} inverse sampling)?
  \textbf{A.} (1) Sample ν\_t\textbar ν\_u from the noncentral χ². (2)
  Sample ∫\_u\^{}t ν\_s ds given (ν\_u,ν\_t) --- the hard step. (3)
  Recover ∫√ν dW\^{}ν=(1/γ)(ν\_t−ν\_u−κ(t−u)+κ∫ν ds) from the CIR
  integral. (4) Sample ∫√ν dW\^{}S\textbar∫ν ds \textasciitilde{} N(0,∫ν
  ds), then assemble S\_t. Step 2 is solved via the known closed-form cf
  of the integrated variance (a ratio of modified Bessel functions);
  invert it to the CDF F --- originally by trapezoidal Fourier
  inversion, \emph{today by COS} --- and sample by inversion ∫ν
  ds=F\^{}\{−1\}(U), U\textasciitilde Uniform(0,1).
\item
  Name the three tools behind the cf of the integrated variance
  (infinitesimal generator, time change, Girsanov) and say what
  discrepancy each one fixes relative to the squared Bessel process.
  \textbf{A.} (i) Infinitesimal generator --- the CIR generator
  κ(−x)f'+½γ²xf'\,' is matched against the squared-Bessel generator
  nf'+2xf'\,', whose Laplace transform of the time-integral is known
  (Pitman--Yor). (ii) Time change s=γ²t/4 rescales the
  \emph{second-order} coefficient ½γ²x onto the squared-Bessel value 2x.
  (iii) Girsanov change of measure absorbs the remaining \emph{drift}
  discrepancy (κ(−x) vs n), expressing the CIR integrated-variance
  transform as a ratio of squared-Bessel transforms --- giving the
  Bessel-ratio formula for the cf.
\item
  Give the qualitative effect on the implied-vol surface of each
  parameter: gamma, rho, kappa, nu0, nu-bar. \textbf{A.} γ (vol-of-vol):
  larger γ deepens smile \emph{curvature}. ρ: more negative ρ steepens
  the \emph{skew} (heavier left tail). κ: small effect on smile shape;
  controls how fast ATM vol converges to √ in maturity. ν\_0: dominates
  \emph{short}-maturity level. : dominates \emph{long}-maturity level
  (both mediated by κ). Overall Heston gives heavier log-return tails
  than BS and reproduces the smile.
\end{enumerate}

\hypertarget{lecture-08}{%
\subsection{Lecture 08}\label{lecture-08}}

\begin{enumerate}
\def\labelenumi{\arabic{enumi}.}
\tightlist
\item
  Define the Poisson process N\_t with intensity lambda. State its three
  defining properties and give P(N\_t=k). What are its mean and
  variance? \textbf{A.} N\_t counts events with intensity λ: (i) N\_0=0;
  (ii) independent, stationary increments; (iii)
  N\_t−N\_s\textasciitilde Poisson(λ(t−s)). So P(N\_t=k)=(λt)\^{}k
  e\^{}\{−λt\}/k!, with mean=variance=λt.
\item
  Show that mean = variance = lambda for the Poisson law. Do the
  E{[}X{]} sum by hand. \textbf{A.} E{[}X{]}=Σ\_\{k≥0\} k λ\^{}k
  e\textsuperscript{\{−λ\}/k!=λe}\{−λ\}Σ\_\{k≥1\}λ\textsuperscript{\{k−1\}/(k−1)!=λe}\{−λ\}e\^{}λ=λ.
  Computing E{[}X(X−1){]}=λ² similarly gives
  Var(X)=E{[}X²{]}−E{[}X{]}²=(λ²+λ)−λ²=λ. So mean=variance=λ.
\item
  What is the compensated Poisson process, and why is it (not N\_t
  itself) a martingale? Verify the martingale property explicitly.
  \textbf{A.} N\_t has positive drift (E{[}N\_t{]}=λt), so it is not a
  martingale. The compensated process Ñ\_t=N\_t−λt is. Check:
  E{[}Ñ\_\{t+Δ\}\textbar F\_t{]}=Ñ\_t+E{[}N\_\{t+Δ\}−N\_t{]}−λΔ=Ñ\_t+λΔ−λΔ=Ñ\_t
  (increment independent of F\_t, mean λΔ), and integrability holds. The
  compensator λt is the predictable finite-variation part removed to
  leave a martingale --- the same idea as the risk-neutral jump drift
  correction.
\item
  Distinguish homogeneous, inhomogeneous (time-dependent lambda(t)), and
  Cox/doubly-stochastic Poisson processes. What is the compensator in
  each case? \textbf{A.} Homogeneous: constant λ, compensator λt.
  Inhomogeneous: deterministic λ(t), counts Poisson with integrated
  intensity Λ(a,b)=∫\_a\^{}b λ(s)ds, compensator ∫\_0\^{}t λ(s)ds.
  Cox/doubly-stochastic: λ(t) itself a stochastic process (conditionally
  inhomogeneous given its path), compensator ∫\_0\^{}t λ(s)ds. In all
  three N\_t−∫\_0\^{}t λ(s)ds is a martingale; the Cox case is the AJD
  slot λ(x)=l\_0+l\_1·x where intensity responds to the state.
\item
  Derive the characteristic function of a compound Poisson process Q\_t
  = sum\_\{k=1\}\^{}\{N\_t\} Y\_k by conditioning on N\_t. You should
  reach exp(t\emph{lambda}(nu\_hat(u)-1)). \textbf{A.} Condition on the
  count: E{[}e\^{}\{iuQ\_t\}{]}=Σ\_n
  P(N\_t=n)E{[}e\^{}\{iuΣ\_\{k≤n\}Y\_k\}{]}=Σ\_n ((λt)\^{}n
  e\textsuperscript{\{−λt\}/n!)(u)}n (Y\_k i.i.d.,
  (u)=E{[}e\^{}\{iuY\}{]}). Pull out e\^{}\{−λt\} and recognise the
  exponential series:
  =e\textsuperscript{\{−λt\}e}\{λt(u)\}=exp(tλ((u)−1)).
\item
  Identify the characteristic exponent of the compound Poisson process
  and connect it to the Levy-Khintchine representation. Why does the
  ``-1'' appear, and what condition does nu\_hat(0)=1 enforce?
  \textbf{A.} The exponent is linear in t, so the characteristic
  exponent is ψ(u)=λ((u)−1)=λ∫(e\^{}\{iuy\}−1)dν(y) --- the
  finite-activity special case of Lévy--Khintchine
  ψ(u)=∫(e\^{}\{iuy\}−1)Π(dy) with Lévy measure Π=λν. The ``−1'' is the
  compensation that makes ψ(0)=0, so φ\_\{Q\_t\}(0)=1 as any cf must
  satisfy ((0)=1 just says ν is a probability measure).
\item
  Write the Merton jump-diffusion SDE for S\_t and derive the log-price
  SDE via Ito-for-jumps. Why is there no (1/2)g'`J\^{}2 term in the jump
  part, and why does ln S\_t - ln S\_\{t\^{}-\} = ln(1+J\_t)?
  \textbf{A.} dS\_t=μS\_\{t\textsuperscript{-\}dt+σS\_\{t}-\}dW\_t+J\_t
  S\_\{t\^{}-\}dN\_t, so at a jump S\_t=S\_\{t\^{}-\}(1+J\_t).
  Itô-for-jumps on g=ln: d ln S=(μ−½σ²)dt+σdW+{[}ln S\_t−ln
  S\_\{t\^{}-\}{]}dN\_t. There is no ½g'`J² term because a jump is a
  \emph{finite} move --- we use the exact difference
  g(Y\_\{t\textsuperscript{-\}+J)−g(Y\_\{t}-\}), not a second-order
  Taylor expansion (the ½g'' term only handles the \emph{continuous}
  Brownian part). And ln S\_t−ln
  S\_\{t\textsuperscript{-\}=ln(S\_\{t}-\}(1+J\_t))−ln
  S\_\{t\^{}-\}=ln(1+J\_t).
\item
  Derive the risk-neutral drift correction mu = r - lambda*E{[}J\_t{]}
  from the requirement that e\^{}\{-rt\}S\_t is a Q-martingale.
  \textbf{A.}
  d(e\textsuperscript{\{−rt\}S\_t)=e}\{−rt\}{[}(μ−r)S\_\{t\textsuperscript{-\}dt+σS\_\{t}-\}dW+J\_t
  S\_\{t\^{}-\}dN\_t{]}. The dW term is a martingale; taking E\^{}Q and
  using E{[}dN\_t{]}=λdt and J⊥N gives drift
  e\textsuperscript{\{−rt\}S\_\{t}-\}{[}(μ−r)+λE\^{}Q{[}J{]}{]}dt.
  Setting it to zero: μ=r−λE\^{}Q{[}J\_t{]}. The jump compensator
  λE{[}J{]} is the mean rate at which jumps push the price up;
  subtracting it keeps the expected return at r.
\item
  Common trap: in the Merton drift correction, is it E{[}J{]} or
  E{[}ln(1+J){]} that gets compensated? If
  ln(1+J)\textasciitilde N(mu\_J,sigma\_J\^{}2), what is kappa =
  E{[}J{]}? \textbf{A.} It is E{[}J{]} (the arithmetic multiplier
  J=S/S\_\{t\^{}-\}−1), \emph{not} E{[}ln(1+J){]}. With
  ln(1+J)=ξ\textasciitilde N(μ\_J,σ\_J²), 1+J=e\^{}ξ, so
  κ=E{[}J{]}=E{[}e\^{}ξ{]}−1=e\^{}\{μ\_J+½σ\_J²\}−1. The compensator is
  λκ=λ(e\^{}\{μ\_J+½σ\_J²\}−1). Confusing E{[}J{]} with E{[}ξ{]}=μ\_J is
  the classic exponent error.
\item
  Write Merton's lognormal-jump cf of ln S\_T and Kou's
  double-exponential cf.~State Kou's nu\_hat(u) and explain the roles of
  p,q,eta\_1,eta\_2 and why eta\_1\textgreater1 is required. \textbf{A.}
  Merton ((u)=e\^{}\{iuμ\_J−½σ\_J²u²\}, κ=e\^{}\{μ\_J+½σ\_J²\}−1):
  φ(u)=exp(iu(x+(r−½σ²−λκ)τ)−½σ²u²τ+λτ(e\^{}\{iuμ\_J−½σ\_J²u²\}−1)).
  Kou: up-jumps Exp(η\_1) w.p. p, down-jumps −Exp(η\_2) w.p. q=1−p, so
  (u)=p η\_1/(η\_1−iu)+q η\_2/(η\_2+iu), and
  φ(u)=exp(iu(x+(r−½σ²−λκ)τ)−½σ²u²τ+λτ((u)−1)). η\_1 governs the
  up-tail decay, η\_2 the down-tail, p/q the up/down probabilities;
  η\_1\textgreater1 is required so E{[}e\^{}X{]}=E{[}1+J{]} is finite
  (the jump multiplier has finite mean, κ finite).
\item
  Compare Merton vs Kou: tail shape, symmetry, and the decay rate of the
  cf in u. Why does the cf decay rate matter for COS? \textbf{A.}
  Merton: Gaussian (symmetric, smooth) jumps; cf decays
  super-exponentially like e\^{}\{−cu²\}. Kou: exponential tails
  (heavier, asymmetric), and the cf decays only rationally like 1/u².
  The cf decay rate sets the decay of the COS cosine coefficients, hence
  the number of terms N for a given accuracy: Merton needs few terms,
  Kou needs more; Kou also fattens the tails so the cumulant range
  {[}a,b{]} is wider.
\item
  How does the jump-diffusion cf enter the COS formula? At which
  arguments is phi evaluated, and what is the ONLY model-specific input
  COS needs? \textbf{A.} The cf is the only model-specific input:
  v≈e\^{}\{−rτ\}Σ' Re\{φ(kπ/(b−a))e\^{}\{−ikπa/(b−a)\}\}V\_k, with φ
  evaluated at the grid arguments u\_k=kπ/(b−a) (and spot carried as
  φ(u)e\^{}\{iux\_0\}). Nothing else about the jump structure is needed
  --- supply Merton/Kou φ and the European COS machinery is unchanged
  (slide 32: COS works as long as the cf is known).
\item
  Short-maturity COS trap: write c\_1, c\_2, c\_4 for Merton's
  log-price. Why can a range {[}a,b{]} built from c\_2 alone be too
  narrow at small tau, and how does the jump c\_4 fix it? (Connect to
  the Assignment 2 bug.) \textbf{A.} Merton cumulants:
  c\_1=(r−½σ²−λκ)τ+λτμ\_J, c\_2=σ²τ+λτ(μ\_J²+σ\_J²),
  c\_4=λτ(μ\_J⁴+6μ\_J²σ\_J²+3σ\_J⁴). At small τ the diffusion part of
  c\_2 (=σ²τ) shrinks, but jumps can still dominate the tails; a range
  built from c\_2 alone (forgetting the jump kurtosis c\_4, or taking L
  too small) is too narrow, so the truncation-error floor stays high and
  the price is wrong \emph{no matter how many terms N} --- exactly the
  Assignment-2 bug. Keeping the √c\_4 term (whose pure-jump contribution
  fattens the tails) and a generous L widens {[}a,b{]} correctly.
\item
  State the PIDE for a jump-diffusion option value. Which term is
  non-local, and why does this motivate Fourier methods over a PIDE grid
  solve? \textbf{A.} V\_t+(r−λE{[}J{]})S
  V\_S+½σ²S²V\_SS−rV+λE\^{}Q{[}V(S(1+J),t)−V(S,t){]}=0. The last term
  --- an integral of V against the jump law --- is \emph{non-local}: it
  couples the value at S to values at all S(1+J). On a grid this is a
  dense coupling (expensive); a transform method instead reads φ once
  and integrates against the payoff, so Fourier/COS sidesteps the
  non-local PIDE entirely.
\item
  Define an affine process and write the AJD SDE with all four affine
  conditions (mu, sigma sigma\^{}T, lambda(x), R(x)). What is the jump
  transform theta(c)? \textbf{A.} Affine: the conditional cf is
  exponential-affine,
  E{[}e\^{}\{u·X\_T\}\textbar X\_t=x{]}=e\^{}\{α+β·x\}. The AJD is
  dX=μ(X)dt+σ(X)dW+dZ with all coefficients affine: drift
  μ(x)=K\_0+K\_1x; covariance
  (σσᵀ)\emph{\{ij\}=(H\_0)}\{ij\}+(H\_1)\_\{ij\}·x; jump intensity
  λ(x)=l\_0+l\_1·x (with jumps J\textasciitilde ν); short rate
  R(x)=ρ\_0+ρ\_1·x. The jump law enters only through the jump transform
  θ(c)=∫e\^{}\{c·z\}dν(z) (the  with complex argument c=β).
\item
  State the AJD result phi = exp(alpha + beta.x) and the Riccati ODEs
  for (alpha,beta) with boundary conditions. Which term is quadratic,
  and where does (theta(beta)-1) come from? \textbf{A.} φ=e\^{}\{α+β·x\}
  with =ρ\_1−K\_1ᵀβ−½βᵀH\_1β−l\_1(θ(β)−1),
  =ρ\_0−K\_0·β−½βᵀH\_0β−l\_0(θ(β)−1), boundary β(T,T)=u, α(T,T)=0. The
  β equation is a Riccati ODE --- the \emph{quadratic} term is the
  ½βᵀH\_1β diffusion term. The (θ(β)−1) comes from the expected jump
  contribution: E{[}(φ(X+J)−φ(X))dN{]}=φ(θ(β)−1)λ(x)dt, the
  compensator-type term, identical to the (−1) of the compound-Poisson
  cf.
\item
  Recover the Black-Scholes log-price cf by solving the AJD ODEs with no
  jumps (X=ln S). Show beta=iu and the alpha you obtain. \textbf{A.} No
  jumps (l\_0=l\_1=0), X=ln S. The ODEs give =0 with β(0)=iu, so β≡iu,
  and =(r−½σ²)β+½σ²β²−r. With β=iu: α(u,τ)=(r−½σ²)iuτ−½σ²u²τ−rτ. Hence
  φ(u,τ)=exp((r−½σ²)iuτ−½σ²u²τ−rτ+iuX) --- the GBM log-price cf,
  recovered as the degenerate AJD. Adding l\_0(θ(iu)−1) reproduces the
  Merton cf.
\item
  The (theta-1) / compensator motif appears at three scales in this
  lecture (compensated Poisson, drift correction, AJD ODE). State all
  three and explain why they are ``the same idea''. \textbf{A.} (i)
  Compensated Poisson: N\_t−λt. (ii) Risk-neutral drift correction:
  μ=r−λE{[}J{]}. (iii) AJD ODEs: the −l(θ(β)−1) term. All three subtract
  the \emph{predictable mean rate of the jumps} (the compensator) so
  that the relevant object --- the counting process, the discounted
  stock, the transform's drift --- becomes a martingale. They are the
  single idea ``remove the predictable jump drift'' at three scales.
\end{enumerate}

\hypertarget{cluster-5-american-options}{%
\section{Cluster 5 --- American
options}\label{cluster-5-american-options}}

\hypertarget{lecture-09}{%
\subsection{Lecture 09}\label{lecture-09}}

\begin{enumerate}
\def\labelenumi{\arabic{enumi}.}
\tightlist
\item
  Define the American early-exercise problem as an optimal-stopping
  problem. Write v(S,t)=max(Λ, continuation) and v\_0 = sup\_τ
  E\^{}Q{[}e\^{}\{-rτ\}Λ(S\_τ,τ){]}. What makes τ* ``free''? \textbf{A.}
  At each instant the holder takes the larger of exercising now and
  holding: v(S,t)=max\{Λ(S,t), continuation\}, where continuation is the
  discounted risk-neutral expectation of v an instant later. Globally
  v\_0=sup\_\{τ∈T\_\{{[}0,T{]}\}\}
  E\^{}Q{[}e\^{}\{−rτ\}Λ(S\_τ,τ)\textbar S\_0{]}, the supremum over all
  stopping times. τ* is ``free'' because the optimal stopping rule is
  not given in advance --- it is itself an output of the problem,
  determined by where continuation falls to intrinsic.
\item
  What is the free boundary S\emph{(t)? Characterise the continuation
  region (v\textgreater Λ, BS PDE holds) vs the stopping region (v=Λ).
  State the smooth-pasting conditions at S}(t) for the put. \textbf{A.}
  S\emph{(t) is the optimal exercise boundary, unknown a priori (hence
  ``free''), separating the continuation region C (hold,
  v\textgreater Λ, where v satisfies the BS PDE with equality) from the
  stopping region S (exercise, v=Λ). For the put:
  S\textless S}(t)⇒exercise, v=K−S; S\textgreater S\emph{(t)⇒hold,
  v\textgreater K−S. At the boundary value and slope match continuously
  --- the smooth-pasting conditions v(S},t)=K−S*, ∂\_S v(S*,t)=−1.
\item
  PROVE that an American call on a non-dividend stock equals the
  European call. Give BOTH the slide dominating-strategy argument and
  the parity lower-bound C\^{}eur = S - Ke\^{}\{-r(T-t)\} + P\^{}eur
  \textgreater{} S - K. \textbf{A.} Dominating strategy: instead of
  exercising at t for S\_t−K, short the stock (receive S\_t) and keep
  the call; at T close the short paying min(K,S\_T)≤K. The alternative
  gives S\_t now and ≤K later, which strictly dominates S\_t−K now once
  interest on the deferred K is counted --- so early exercise is
  dominated at every t\textless T, hence worthless. Lower-bound version:
  by parity and P\^{}eur≥0,
  C\textsuperscript{eur=S−Ke}\{−r(T−t)\}+P\textsuperscript{eur≥S−Ke}\{−r(T−t)\}\textgreater S−K
  for r\textgreater0. The live call exceeds its intrinsic, so the
  early-exercise right is never used: C\textsuperscript{am=C}eur.
\item
  Why does the same argument FAIL for the put, so early exercise can be
  optimal? Use parity P\^{}eur = Ke\^{}\{-r(T-t)\} - S + C\^{}eur and
  Ke\^{}\{-r(T-t)\} \textless{} K. What is the economic intuition?
  \textbf{A.} Parity gives
  P\textsuperscript{eur=Ke}\{−r(T−t)\}−S+C\^{}eur, and
  Ke\^{}\{−r(T−t)\}\textless K, so the live put can be worth \emph{less}
  than its intrinsic K−S. When that happens, exercising now and banking
  the full K to earn interest beats holding --- early exercise is
  strictly optimal, so there is no European--American equivalence.
  Intuition: the put holder \emph{wants the strike cash K early} to earn
  interest, whereas the call holder wants to \emph{defer paying K} ---
  the asymmetry is entirely the sign of the strike cash flow.
\item
  Why does delta-hedging make the European BS relation an EQUALITY but
  the American put only an INEQUALITY? Walk through both arbitrage
  directions and explain why only one survives. \textbf{A.} For a
  European option the delta-hedged position is perfectly riskless, so it
  must earn r exactly --- equality. For the American put, hedging does
  not remove the holder's early-exercise right. If the hedged option
  \emph{out-earns} cash, an arbitrageur longs it (owning the exercise
  right) and locks the surplus --- a genuine arbitrage, competed away.
  If it \emph{under-earns} cash, arbitraging would require
  \emph{shorting} the put, but that hands the exercise decision to the
  counterparty, who can exercise during {[}t,t+Δt{]} --- no riskless
  profit. Only the second case survives, giving LP\_Am≤0, hence an
  inequality.
\item
  Write the LCP for the American put: (v-Λ)≥0, Lv≤0, (v-Λ)·Lv=0.
  Interpret each line, why it equals an obstacle problem, and state
  terminal/boundary conditions. \textbf{A.} (v−Λ)≥0 (value dominates
  payoff, else buy-and-exercise arbitrage), Lv≤0 (the BS operator is a
  supersolution, from the hedging inequality), (v−Λ)·Lv=0
  (complementarity: at each point one holds with equality). In C,
  v\textgreater Λ ⇒ Lv=0 (BS PDE holds); in S, v=Λ ⇒ Lv≤0 (payoff is not
  a BS solution, exercise optimal). It is an obstacle problem because v
  is the smallest supersolution of L lying above the obstacle Λ.
  Conditions: terminal v(S,T)=max(K−S,0); boundary v(0,t)=K, v(∞,t)=0.
\item
  State the Bermudan backward recursion in log-space:
  c(x,t\_\{n-1\})=e\^{}\{-rΔt\}E{[}v(X(t\_n))\textbar x{]}, v=max(Λ,c).
  What single quantity do all four methods compute differently?
  \textbf{A.} With X(t\_n)=ln(S(t\_n)/K), Λ(x,t)=max(αK(e\^{}x−1),0),
  the recursion for n=N,\ldots,2 is
  c(x,t\_\{n−1\})=e\^{}\{−rΔt\}E{[}v(X(t\_n),t\_n)\textbar X(t\_\{n−1\})=x{]},
  v(x,t\_\{n−1\})=max(Λ(x,t\_\{n−1\}),c(x,t\_\{n−1\})). All four methods
  (tree, FD/LCP, LSM, COS) run exactly this Bellman recursion; they
  differ \emph{only} in how they compute the conditional expectation c
  --- the continuation value.
\item
  Binomial tree: write the backward step with max(intrinsic, discounted
  continuation). What is the ONLY line that differs from a European
  tree? How does p=(e\^{}\{rΔt\}-d)/(u-d) arise? \textbf{A.}
  V\_m\textsuperscript{n=max(K−S\_m}n,
  e\^{}\{−rΔt\}{[}pV\_\{m+1\}\textsuperscript{\{n+1\}+(1−p)V\_m}\{n+1\}{]}).
  The conditional expectation is the two-point binomial average; the
  \emph{only} difference from a European tree is the outer
  max(intrinsic, ·) inserted at every node. p=(e\^{}\{rΔt\}−d)/(u−d)
  arises from the risk-neutral / martingale condition
  E\textsuperscript{Q{[}S\_\{n+1\}\textbar S\_n{]}=e}\{rΔt\}S\_n,
  i.e.~pu+(1−p)d=e\^{}\{rΔt\}.
\item
  Finite-difference LCP: write the projected update. Why can't you just
  solve the implicit linear system and clip once? Explain PSOR ---
  projecting INSIDE each SOR sweep. \textbf{A.} Projected update:
  V\_m\textsuperscript{n=max(K−S\_m}n,
  g(V\_\{m+1\}\textsuperscript{\{n+1\},V\_m}\{n+1\},V\_\{m−1\}\^{}\{n+1\}))
  with g the implicit BS stencil --- the discrete LCP. For an
  \emph{implicit} scheme the node m couples to its neighbours at the
  \emph{same} time level through a tridiagonal system, and the
  constraint V≥Λ must hold simultaneously with that solve; solving the
  unconstrained system and clipping once gives the wrong answer because
  the clipped values feed back into the coupling. PSOR (Projected SOR)
  runs SOR sweeps but, after each node's update, immediately projects
  V\_m←max(Λ\_m, V\_m\^{}\{SOR\}); projecting \emph{inside} the
  iteration converges to the true LCP solution.
\item
  LSM mechanics: explain the backward regression of discounted future
  cash flows on basis functions. Why regress only on ITM paths? Why is
  the per-path decision Λ\textgreater ĉ (not the fitted value) the cash
  flow? \textbf{A.} Working backward, at each date regress each path's
  realised discounted future cash flow Y (from following the
  already-determined later policy, discounted to t\_n) on basis
  functions of the current state X=S\_\{t\_n\} (e.g.~1, X, X²); the fit
  ĉ(x)=Σâ\_j L\_j(x) estimates the continuation value. Regress only on
  ITM paths because OTM paths have zero intrinsic and would never
  exercise --- restricting fits the conditional expectation only where
  the decision matters (cheaper and more accurate). The realised
  \emph{cash flow} used is the path's actual intrinsic when
  Λ\textgreater ĉ, not the fitted ĉ: ĉ is only a noisy decision rule,
  while the realised payoff is the unbiased value collected on that
  path.
\item
  LSM 8-path example: how is Y constructed (which discount factor, which
  future date)? Why does the t=1 regression depend on the t=2 exercise
  decisions? Interpret 0.1144 vs European 0.0564. \textbf{A.} At t=2, Y
  is the t=3 payoff discounted one period (×e\^{}\{−0.06\}=0.94176),
  X=S\_\{t=2\}; fit gives exercise on paths 4,6,7. At t=1, Y is the
  \emph{subsequent} realised cash flow --- which by the t=2 decisions
  occurs at t=2 or t=3 --- discounted back to t=1; so the t=1 regression
  depends on the t=2 exercise decisions because they determine when each
  path's future cash flow actually lands. Final: discount each path's
  first-exercise cash flow to t=0 and average → American put 0.1144, vs
  European 0.0564 (only the t=3 payoffs). The American is
  \textasciitilde twice the European: that gap is the entire value of
  the early-exercise right.
\item
  Is LSM biased? Explain why a suboptimal finite-basis policy gives a
  LOWER bound, and the competing in-sample foresight upward bias. How do
  you get a clean out-of-sample estimate? \textbf{A.} A finite basis
  cannot represent the true continuation surface exactly, so the
  estimated stopping policy is suboptimal and a suboptimal rule can only
  \emph{lose} value --- giving a lower bound. There is a competing
  \emph{upward} (foresight) bias if the price is read off the same cash
  flows used to fit the regression (in-sample look-ahead). Clean
  practice: use the regression only for the exercise \emph{decision} and
  value the realised cash flows along each path's stopping time; for a
  strictly out-of-sample estimate, freeze the policy and re-simulate
  fresh paths.
\item
  COS Bermudan: derive c(x,t\_\{n-1\}) ≈ e\^{}\{-rΔt\} Σ'
  Re\{φ(kπ/(b-a)) e\^{}\{ikπ(x-a)/(b-a)\}\} V\_k(t\_n). What ARE the
  V\_k(t\_n) --- relate to the Lecture-6 European V\_k. \textbf{A.}
  Write c(x,t\_\{n−1\})=e\^{}\{−rΔt\}∫v(y,t\_n)f(y\textbar x)dy --- the
  European COS integral with the payoff replaced by the next-step value
  v(·,t\_n). Inserting the cosine expansion of f (coefficients read off
  φ) gives ĉ(x,t\_\{n−1\})=e\^{}\{−rΔt\}Σ'\_\{k\}
  Re\{φ(kπ/(b−a))e\^{}\{ikπ(x−a)/(b−a)\}\}V\_k(t\_n). Here
  V\_k(t\_n)=(2/(b−a))∫\_a\^{}b v(y,t\_n)cos(kπ(y−a)/(b−a))dy are the
  cosine coefficients of the option \emph{value} at t\_n.~At t\_N they
  are exactly the Lecture-6 European payoff coefficients; for
  n\textless N they are built recursively --- so it is a backward
  recursion on the coefficient vector, not on a price grid.
\item
  COS early-exercise point: define x\_n* by
  c(x\_n\emph{,t\_n)=Λ(x\_n},t\_n). Write the split
  V\_k(t\_n)=G\_k(a,x\_n\emph{)+C\_k(x\_n},b,t\_n) for a put. Which
  piece is closed-form (χ\_k,ψ\_k), which needs the previous
  coefficients? How is x\_n* found? \textbf{A.} x\_n* solves
  c(x\_n\emph{,t\_n)=Λ(x\_n},t\_n) --- the discrete avatar of the free
  boundary S\emph{(t). For a put, exercise is on the low-y side:
  V\_k(t\_n)=G\_k(a,x\_n})+C\_k(x\_n\emph{,b,t\_n), where G\_k is the
  cosine integral of the payoff Λ=αK(e\^{}y−1) and C\_k that of the
  continuation c.~G\_k is closed-form --- a combination of the Lecture-6
  χ\_k (the e\^{}y cos integral) and ψ\_k (the cos integral) on the cut
  interval. C\_k needs the }previous* step's coefficients
  V\_l(t\_\{n+1\}) (it is built from ĉ). x\_n* is found by Newton's
  method (one scalar root-find; the derivative of ĉ is the
  differentiated cosine sum).
\item
  COS algorithm: initialisation, backward loop (Newton for x\_n\emph{,
  FFT for C\_k), reconstruction. Why is the cost O(NK log₂K) not O(NK²)?
  What gives the FFT (Toeplitz/Hankel kernel)? \textbf{A.} Initialise at
  t\_N with the pure payoff coefficients V\_k(t\_N)=G\_k (Lecture-6
  European). Backward loop n=N−1,\ldots,1: (a) find x\_n} by Newton from
  ĉ=Λ; (b) form V\_k(t\_n) via the split --- G\_k closed-form, C\_k by
  the FFT-accelerated product on \{V\_l(t\_\{n+1\})\}. Reconstruct:
  v(x\_0,t\_0)=e\^{}\{−rΔt\}Σ'
  Re\{φ(kπ/(b−a))e\^{}\{ikπ(x\_0−a)/(b−a)\}\}V\_k(t\_1). The C\_k
  coupling matrix M\_\{k,l\} is Hankel-plus-Toeplitz (depends on l±k),
  so the matrix--vector product is a discrete convolution done by FFT in
  O(K log₂K) instead of the naive O(K²); over N steps the cost is O(NK
  log₂K).
\item
  Why does the cumulant range {[}a,b{]}=c1±L√(c2+√c4) caveat carry into
  the Bermudan COS recursion and bite harder per-step? Tie to the
  Assignment-2 bug: c2=σ²Δt→0, c4=0, no floor; adding K terms does NOT
  help. \textbf{A.} The range is still set from the cumulants of X\_T,
  but now it must contain the support over the \emph{whole} backward
  recursion and the Newton x\_n* must land inside {[}a,b{]}. At short
  maturity (or small per-step Δt) c\_2=σ²Δt→0, so the box collapses
  toward c\_1; for a Gaussian model c\_4=0 means no √c\_4 floor --- the
  same failure mode as the Assignment-2 bug. Adding cosine terms K
  cannot cure a range (truncation) error; the fix is a wider L or a
  minimum half-width floor. It bites per-step because each of the N
  steps re-uses the same possibly-too-narrow box.
\item
  Richardson extrapolation Bermudan→American: why does every method here
  price a Bermudan? Write the 4-point scheme and explain what it
  cancels. \textbf{A.} Every method discretises exercise to a finite
  monitoring set, so each prices a \emph{Bermudan}; the American is the
  N→∞ limit. Rather than push N large, take Bermudan prices v(M) at
  M=2\textsuperscript{d,2}\{d+1\},2\textsuperscript{\{d+2\},2}\{d+3\}
  and combine:
  v\_Am(d)=(1/21)(64v(2\textsuperscript{\{d+3\})−56v(2}\{d+2\})+14v(2\textsuperscript{\{d+1\})−v(2}d)).
  The weighted combination cancels the leading error terms in the
  monitoring spacing, accelerating convergence to the
  continuous-American limit from a handful of cheap fast Bermudan
  computations.
\end{enumerate}

\hypertarget{cluster-6-exotics}{%
\section{Cluster 6 --- Exotics}\label{cluster-6-exotics}}

\hypertarget{lecture-11}{%
\subsection{Lecture 11}\label{lecture-11}}

\begin{enumerate}
\def\labelenumi{\arabic{enumi}.}
\tightlist
\item
  Give the precise payoff of an arithmetic-average Asian call and a
  geometric-average Asian call. \textbf{A.} Fixed-strike
  arithmetic-average call: ((1/n)Σ\_\{i=1\}\^{}n S\_\{t\_i\}−K)\^{}+.
  Geometric-average call: ((∏\emph{\{i=1\}\^{}n
  S}\{t\_i\})\textsuperscript{\{1/n\}−K)}+. (Floating-strike variants
  replace K by a multiple of S\_T.) Both are path-dependent on the
  monitoring dates t\_1,\ldots,t\_n.
\item
  What economic need does an Asian option serve that a vanilla cannot,
  and why is the premium lower? (airline fuel hedge; average exposure;
  effective variance reduced) \textbf{A.} It matches an exposure to an
  \emph{average} rather than a terminal price --- e.g.~an airline
  hedging a month of fuel procurement buys one Asian call on the monthly
  average instead of \textasciitilde20 daily vanillas, matching its real
  cost (1/n)ΣS\_\{t\_i\} and reducing fixing-manipulation risk. The
  premium is lower because the average is less volatile than the
  terminal price (continuous-monitoring variance σ²T/3, one-third of the
  vanilla), so the convex payoff is worth less.
\item
  Derive the geometric-Asian closed form: why is G lognormal, and what
  are mu\_G and sigma\_G\^{}2 for equally spaced dates? Show
  sigma\_G\^{}2 -\textgreater{} sigma\^{}2 T/3 in the continuous limit.
  \textbf{A.} ln G=(1/n)Σ{[}(r−q−½σ²)t\_i+σW\_\{t\_i\}{]} is an affine
  function of the jointly Gaussian (W\_\{t\_1\},\ldots,W\_\{t\_n\}),
  hence normal --- so G is lognormal and prices by a Black-style
  formula. μ\_G=ln S\_0+(r−q−½σ²), σ\_G²=(σ²/n²)Σ\_iΣ\_j
  min(t\_i,t\_j). For equally spaced t\_i=iΔt: =(n+1)T/2n,
  σ\_G²=σ²Δt(n+1)(2n+1)/6n² → σ²T/3 as n→∞. The σ²T/3 (vs σ²T) is
  exactly why the Asian is cheaper, and the geometric Asian is the
  standard MC control variate for the arithmetic one (Kemna--Vorst).
\item
  Why does the arithmetic Asian have NO closed form / no closed-form CF,
  and what does that imply for method choice? (sum of lognormals not
  lognormal) \textbf{A.} A\_T=(1/n)ΣS\_\{t\_i\} is a \emph{sum} of
  lognormals, which is not lognormal and has no elementary density ---
  and crucially no closed-form characteristic function even under affine
  models. So one cannot price it by Black or apply COS directly to A\_T;
  the options are moment-matching, the (Vecer-reduced) PDE, COS with a
  cumulant-approximated/recursive cf, or Monte Carlo.
\item
  List the four numerical methods for arithmetic Asians and the one-line
  reason for each; what is the standard control variate for the MC?
  (geometric Asian) \textbf{A.} (i) Moment matching
  (Turnbull--Wakeman/Levy) --- fit a lognormal to the first two moments
  of A\_T; fast but an approximation that degrades in the tails. (ii)
  PDE --- 2-D in (S,A), reduced exactly to 1-D by Vecer's similarity
  trick. (iii) COS/Fourier (Zhang--Oosterlee) --- cumulant-approximated
  φ\_\{A\_T\} or backward recursion over monitoring dates with the
  one-step conditional cf.~(iv) Monte Carlo --- most flexible for path
  dependence, with the \emph{geometric}-Asian closed form as control
  variate (essentially mandatory variance reduction).
\item
  State the Fouque-Han 3D Asian PDE and identify the four blocks. Then
  explain Vecer's reduction: what does q(t) replicate, and how does it
  collapse the path-dependent state I\_t? \textbf{A.} With (S,Y,I),
  I\_t=∫\emph{0\^{}t S\_u du, the floating-strike Asian solves
  ∂\emph{tP+½f²s²∂}\{ss\}P+ρβfs∂}\{sy\}P+½β²∂\_\{yy\}P+r(s∂\_sP−P)+(α(m−y)−βΛ)∂\_yP+s∂\_IP=0,
  terminal (s−I/T)\^{}+. Four blocks: second-order (s,y) diffusion
  (cross-term carries ρ); first-order risk-neutral drift in s and
  risk-adjusted drift in y; the pure-drift term s∂\_IP (since dI=Sdt);
  and −rP discounting. Vecer: hold a \emph{time-deterministic} number of
  shares q(t)=(1−e\^{}\{−r(T−t)\})/(rT) in a self-financing portfolio so
  that X\_T=(1/T)∫S du−K\_2 --- i.e.~q(t) replicates the running
  average. Trading this deterministic q collapses the path-dependent
  state I\_t into the single stochastic wealth X\_t; switching to the
  stock numéraire and using homogeneity reduces the 3-D PDE to 2-D (1-D
  in the constant-vol limit, the Vecer Asian PDE).
\item
  Give both digital payoffs (cash-or-nothing, asset-or-nothing). Derive
  the cash-or-nothing call price = e\^{}\{-rT\} N(d2) from Q(S\_T
  \textgreater{} K). \textbf{A.} Cash-or-nothing call:
  1\_\{S\_T\textgreater K\} (× notional Q). Asset-or-nothing call: S\_T
  1\_\{S\_T\textgreater K\}. The cash-or-nothing pays \$1 on
  \{S\_T\textgreater K\}, so C\_dig=e\^{}\{−rT\}Q(S\_T\textgreater K).
  Under GBM, \{S\_T\textgreater K\}=\{Z\textgreater−d\_2\} with
  d\_2=(ln(S\_0/K)+(r−q−½σ²)T)/(σ√T), so Q(S\_T\textgreater K)=N(d\_2)
  and C\_dig=e\^{}\{−rT\}N(d\_2). A digital prices a risk-neutral
  probability; the asset-or-nothing is S\_0e\^{}\{−qT\}N(d\_1), and
  vanilla = asset-or-nothing − K·cash-or-nothing.
\item
  A digital has a clean closed form. Why is it nonetheless considered a
  hard/dangerous product? (delta spike, gamma sign-flip, pin risk;
  call-spread notional 1/eps -\textgreater{} infinity near K at expiry)
  \textbf{A.} Valuation is trivial (e\^{}\{−rT\}N(d\_2)); the
  \emph{hedge} is not. The replicating call spread has notional 1/ε→∞ as
  the spread tightens, and near expiry near the strike the delta spikes
  (Dirac-like) and gamma flips sign --- ``pin risk''. A desk
  delta-hedging a sold digital must trade a huge, sign-changing position
  in a vanishingly small neighbourhood of K as t→T, impossible to
  execute cleanly; in practice they over-hedge by pricing/hedging a
  (super-replicating) call spread, not the exact digital. The difficulty
  is risk management, not pricing.
\item
  State the Breeden-Litzenberger link dC/dK = -e\^{}\{-rT\} Q(S\_T
  \textgreater{} K) and explain the centred call-spread static
  replication (and why centring gives 2nd-order accuracy). \textbf{A.}
  ∂C/∂K=−e\^{}\{−rT\}Q(S\_T\textgreater K)=−C\_dig: a digital is minus
  the strike-derivative of vanilla prices, linking it to the implied
  (risk-neutral) distribution. Static replication:
  1\_\{S\_T\textgreater K\}≈{[}(S\_T−K+ε/2)\textsuperscript{+−(S\_T−K−ε/2)}+{]}/ε,
  i.e.~buy 1/ε calls at K−ε/2 and sell 1/ε at K+ε/2. \emph{Centring} the
  spread symmetrically around K makes the leading O(ε) error cancel (the
  approximation is a centred difference of the call price in K), giving
  second-order O(ε²) accuracy; a one-sided spread is only first-order.
\item
  Which numerical method is most natural for a digital and why? (COS:
  indicator payoff has closed-form cosine coefficients; it's the
  European COS method as a special case) \textbf{A.} COS. The indicator
  payoff has closed-form cosine coefficients H\_k (an elementary
  cos-integral over the in-the-money half-range), so a digital is just
  the European COS method with a different H\_k --- no discontinuity
  penalty on the payoff side, since the coefficients are computed
  analytically rather than by sampling the jump. (Closed form exists
  under GBM; COS handles any model with a known cf.)
\item
  Give the basket and worst-of payoffs. Show Var(B\_T) = sum\_i sum\_j
  w\_i w\_j Cov(S\_i,S\_j) and explain why correlation is the dominant
  risk. \textbf{A.} Basket call: (Σ\_i w\_i S\_i(T)−K)\^{}+. Worst-of
  call: (min\_i S\_i(T)−K)\^{}+. For B\_T=Σ\_i w\_i S\_i(T),
  Var(B\_T)=Σ\_iΣ\_j w\_i w\_j Cov(S\_i,S\_j)=Σ\_i
  w\_i²Var(S\_i)+Σ\_\{i≠j\}w\_i w\_j ρ\_\{ij\}√(Var S\_i·Var S\_j). The
  off-diagonal ρ\_\{ij\} terms drive the variance: positively correlated
  baskets are more volatile (dearer), weakly/negatively correlated ones
  diversify. A worst-of is even more correlation-sensitive (ρ→1 makes it
  a single-asset option; low ρ drags the minimum down). The 2008
  correlation break-down is what made these books dangerous.
\item
  Why is Monte Carlo (not PDE or COS) the method for a
  \textgreater=4-asset worst-of? State the COS availability honestly.
  (curse of dimensionality; Ruijter-Oosterlee 2D only, no dedicated COS
  for n\textgreater=4) \textbf{A.} PDE and COS both use a tensor-product
  grid whose cost grows like M\^{}n --- the curse of dimensionality
  kills them beyond n≈2--3. Honestly: COS exists for 2-asset
  basket/call-on-min (Ruijter--Oosterlee 2012) and a 3-asset spread
  (Pellegrino--Sabino 2016), but there is \emph{no dedicated COS method
  for n≥4}. MC handles arbitrary dimension and correlation naturally
  (Cholesky-correlate normals, simulate, evaluate (min\_i S\_i−K)\^{}+),
  enhanced with antithetics/QMC/control variates --- hence the industry
  standard for the multi-asset book.
\item
  In a worst-of autocallable, is the embedded option a call-on-min or a
  worst-of put? Explain who is short what. \textbf{A.} A worst-of
  \emph{put}, not a clean call-on-min. The investor is \emph{short}
  downside on the weakest index (e.g.~redemption is reduced if min\_i
  S\_i(T)/S\_i(0) falls below a 60\% barrier), implicitly \emph{selling}
  correlation and downside risk; the high coupon funds that sold risk.
  The issuing bank is long that protection (short correlation in its
  book). Be ready to state that the traded structure is a worst-of put,
  distinct from the textbook call on the min.
\item
  Give the cliquet payoff with local cap C and floor F. Why is it
  forward-skew / forward-volatility sensitive, and why is MC with
  local/stochastic vol the practical method? \textbf{A.} With period
  returns R\_i=S\_\{t\_i\}/S\_\{t\_\{i−1\}\}−1, the cliquet pays Σ\_i
  min(max(R\_i,F),C) (optionally with a global cap/floor on the sum). It
  is essentially a strip of \emph{forward-starting} capped/floored
  options: each l\_i is determined over the future window
  {[}t\_\{i−1\},t\_i{]}, so its value depends on the \emph{forward}
  volatility and forward skew prevailing then, not today's spot-starting
  smile --- an exposure no vanilla portfolio (all spot-starting) can
  replicate. So pricing is highly sensitive to smile \emph{dynamics};
  the practical method is MC with a local/stochastic-vol (or hybrid)
  model that gets the forward skew right, since the reset features
  combine trivially with such models.
\item
  Pre-2008 vs post-2008: which exotics grew, which shrank, and what is
  the unifying post-crisis preference? (hedgeable, modellable,
  explainable, liquid; basket/worst-of and autocallables grew;
  bespoke/long-dated/aggressive cliquets shrank) \textbf{A.} Grew:
  basket/worst-of and autocallables (now dominant in equity structured
  products), driven by retail yield demand outrunning the correlation
  caveats. Shrank: highly bespoke OTC, long-dated path-dependent
  exotics, and aggressive cliquets (now insurance-driven, reduced).
  Barriers/Asians/digitals stayed active. The unifying post-crisis
  preference is for products that are easier to \emph{hedge, model,
  explain to regulators, and more liquid} --- reflecting XVA/funding
  costs, Basel III capital, collateralisation, and the lessons of
  counterparty risk and correlation break-down.
\end{enumerate}

\end{document}
